% ****** Start of file apssamp.tex ******
%
%   This file is part of the APS files in the REVTeX 4.2 distribution.
%   Version 4.2a of REVTeX, December 2014
%
%   Copyright (c) 2014 The American Physical Society.
%
%   See the REVTeX 4 README file for restrictions and more information.
%
% TeX'ing this file requires that you have AMS-LaTeX 2.0 installed
% as well as the rest of the prerequisites for REVTeX 4.2
%
% See the REVTeX 4 README file
% It also requires running BibTeX. The commands are as follows:
%
%  1)  latex apssamp.tex
%  2)  bibtex apssamp
%  3)  latex apssamp.tex
%  4)  latex apssamp.tex
%
\documentclass[%
 reprint,
%superscriptaddress,
%groupedaddress,
%unsortedaddress,
%runinaddress,
%frontmatterverbose, 
%preprint,
%preprintnumbers,
%nofootinbib,
%nobibnotes,
%bibnotes,
 amsmath,amssymb,
 aps,
%pra,
%prb,
%rmp,
%prstab,
%prstper,
%floatfix,
]{revtex4-2}

\usepackage{graphicx}% Include figure files
\usepackage{dcolumn}% Align table columns on decimal point
\usepackage{bm}% bold math
\usepackage{lipsum}
\usepackage{hyperref}% add hypertext capabilities
\usepackage[capitalize]{cleveref}
\crefname{section}{Appendix}{Appendices} 
\usepackage[mathlines]{lineno}% Enable numbering of text and display math
%\linenumbers\relax % Commence numbering lines

%\usepackage[showframe,%Uncomment any one of the following lines to test 
%%scale=0.7, marginratio={1:1, 2:3}, ignoreall,% default settings
%%text={7in,10in},centering,
%%margin=1.5in,
%%total={6.5in,8.75in}, top=1.2in, left=0.9in, includefoot,
%%height=10in,a5paper,hmargin={3cm,0.8in},
%]{geometry}

\begin{document}

\preprint{APS/123-QED}

\title{Non-Thermal Aging and Cooling of Supercooled Liquids in Microcavities}% Force line breaks with \\

\author{Muhammad R. Hasyim, Arianna Damiani, Norah M. Hoffmann}
\email{mh7373@nyu.edu}
\affiliation{The Simons Center for Computational Physical Chemistry, New York University, New York, NY 10003}%


\date{\today}% It is always \today, today,
             %  but any date may be explicitly specified

\begin{abstract}

Supercooled liquids fall out of equilibrium below the glass transition temperature $T_\mathrm{g}$ due to dramatic slowdown in structural relaxation. 
Below $T_\mathrm{g}$, they undergo physical aging---a slow drift of material properties toward equilibrium that can span astonomical timescales, making equilibrium preparation a central challenge in glass physics. 
Here, we show how placing atomistic supercooled liquids inside Fabry-Pérot cavities---optical resonators that confine resonant electomagnetic (EM) modes---induces nonthermal aging via strong-light matter coupling. 
This phenomenon occurs as cavity EM modes selectively pump energy into intramolecular vibrations, creating imbalances that drive intermolecular degrees of freedom into deeper potential minima to maintain equilibrium at all times.
Applying known principles of physical aging, such as time reparameterization softness and material-time translational invariance,  we collapse nonequilibrium correlation functions onto their equilibrium counterparts, revealing that cavity aging traverses sequences of equilibrium states.
Exploiting this, we develop \textit{cavity configurational feedback (\texttt{C2F}) cooling}, where we apply periodic cavity coupling with feedback control on an external reservoir to prepare an equilibrium supercooled liquid at ultralow temperatures ($T \approx  T_\mathrm{g}$).
These results show that strong light-matter interactions provide a transformative route to access ultralow-temperature equilibrium states. %and probe longstanding questions in glass physics.

% The slowdown of relaxation in supercooled liquids limits access to equilibrium at ultralow temperatures, where physical aging dominates. 
% We show that confining glass-forming liquids in optical cavities induces nonthermal aging governed by separated timescales. 
% Classical cavity molecular dynamics (cavMD) reveals that cavity modes drive intramolecular vibrations on femtosecond scales, while thermal baths maintain kinetic equilibrium on picosecond scales. 
% This selective pumping forces intermolecular degrees of freedom into deeper energy minima, effectively cooling the structural subsystem without lowering temperature. 
% The resulting self-quenching mechanism slows relaxation by nearly an order of magnitude, with aging dynamics persisting far beyond equilibrium timescales. Nonequilibrium correlation functions collapse onto stretched exponentials independent of coupling strength, and the Tool–Narayanaswamy model accurately predicts trajectories using only fictive temperatures and equilibrium inputs. 
% Thus, cavity coupling functions as a frequency-selective engineered bath, providing a protocol for configurational cooling into deeply supercooled equilibrium states beyond conventional thermal methods.
% The fields of polaritonics and glassy dynamics have developed independently: the former studies molecular interactions with confined electromagnetic fields, while the latter focuses on the dramatically slow relaxation of supercooled liquids.
% Optical cavities have enabled precise manipulation and engineering of material properties in chemistry and condensed matter, showing promise for engineering relaxation pathways in supercooled liquids.
% Here, we demonstrate that placing a glass-forming mixture of polar molecules within a Fabry-Pérot cavity induces nonthermal aging through selective energy redistribution across various timescales.
% Using classical cavity molecular dynamics, we demonstrate that the cavity channels energy into intramolecular vibrations on femtosecond scales while thermal baths maintain equilibrium on picosecond scales, creating an energy constraint.
% To compensate, intermolecular degrees of freedom access deeper potential energy minima, effectively cooling the structural subsystem.
% This self-quenching mechanism slows structural relaxation by nearly an order of magnitude despite kinetic temperatures remaining at equilibrium values.
% Analysis with the Tool--Narayanaswamy model reveals universal stretched exponential aging, confirming that cavity coupling preserves fundamental glassy characteristics.
% Our findings establish a robust method for controlling relaxation pathways in supercooled liquids through light-matter interactions, with potential applications in preparing deeply supercooled states.
\end{abstract}

%\keywords{Suggested keywords}%Use showkeys class option if keyword
                              %display desired
\maketitle

%\tableofcontents

%\section{Introduction}\label{sec:level1}
%First-level heading:\protect\\ The linebreak was forced \lowercase{via} \textbackslash\textbackslash}
%\section*{Introduction}

\begin{figure*}[t]
    \centering
    \includegraphics[width=\linewidth]{Figure1.pdf}
    \caption{Overview of nonthermal aging and cooling in supercooled liquids. 
    (a) Schematic of a supercooled liquid in a Fabry--Pérot microcavity with mirror spacing $L_\mathrm{c} \sim O(\mu\mathrm{m})$. The cavity mode (frequency $\omega_\mathrm{c}$) couples to molecular vibrations. Both molecular and cavity systems exchange energy with thermal baths at temperature $T$ on timescales $\tau_\mathrm{b}$ and $\tau_\mathrm{c}$, respectively.
    (b) Left: Infrared (IR) absorption spectra $n(\omega)\alpha(\omega)$ at equilibrium ($T = 100$ K) with $\omega_\mathrm{c} = 1560$ cm$^{-1}$ matching the A$_2$ stretch. Increasing coupling $\lambda$ (blue to red) splits the vibrational peak into upper ($\omega_\mathrm{UP}$) and lower ($\omega_\mathrm{LP}$) polariton branches, confirming strong light-matter coupling. Right: Linear scaling of Rabi splitting $\Omega_\mathrm{R}=\omega_\mathrm{UP}-\omega_\mathrm{LP}$ with coupling $\lambda$.
    (c) Evolution of correlation functions comparing out-of-cavity (blue) and in-cavity (orange, $\lambda = 0.141$ a.u.) dynamics during nonthermal aging. Illustrations shows selective pumping of intramolecular vibrations forcing intermolecular degrees of freedom into colder and deeper minima. In cavity, the system gradually ages back to equilibrium where both correlation functions converge.
    (e) The time-evolution of the structural fictive temperature $T_\mathrm{F}(t)$ during a normal quench and the \texttt{C2F} protocol. While a normal quench leads to being trapped in a local minimum, the C2F protocol reaches target at ultrafast timescales.
    }
    \label{fig:fig1}
\end{figure*}

% Controlling relaxation dynamics of supercooled liquids---the metastable state of glass-forming liquids---remains a central challenge in condensed matter physics.
% The challenge stems from the dramatic non-Arrhenius increase in their relaxation times that can span picoseconds to millions of years upon decreasing temperature \cite{angell2000relaxation}. 
% Due to this slowdown, supercooled liquids fall out of equilibrium at the glass transition temperature $T_\mathrm{g}$, limiting access to equilibrium states at ultralow temperatures ($T < T_\mathrm{g}$). 
% At these temperatures, glassy systems display physical aging, where static material properties begin its gradual drift towards the true equilibrium \cite{berthier2011theoretical,biroli2013perspective}.
% While aging studies have revealed fundamental principles such as time reparameterization softness \cite{chamon2002separation,ghimenti2024cleveralgorithmsglassesdo, bohmer2024time}, material-time translational invariance \cite{chamon2002separation,ghimenti2024cleveralgorithmsglassesdo,bohmer2024time}, and violations of the fluctuation-dissipation theorem \cite{chamon2002separation, BarratKob1999, di2000off, grigera1999observation, PhysRevLett.93.160603, wang2006dynamic}, these insights have not resolved the fundamental challenge of preparing equilibrium states themselves.

% Conventional protocols to prepare supercooled liquids require cooling rates that must be sufficiently slow to maintain adiabatic equilibration, yet fast enough to suppress crystallization---competing requirements that impose fundamental limits.
% Only recently have experimental advances, such as physical vapor deposition \cite{ediger2017perspective,rodriguez2022ultrastable}, and computational methods, including swap Monte Carlo algorithms \cite{ninarello2017models,scalliet2022thirty,jung2025numericalinvestigationequilibriumkauzmann}, enabled the preparation of equilibrium states at previously inaccessible temperatures.
% These methods, however, still fall short of reaching the Kauzmann temperature, a hypothetical transition where non-ergodic glass emerges as a true equilibrium phase of matter.
% Thus, protocols to equilibrate supercooled liquids remain the main bottleneck for solving the deepest problems in glass physics. 

Controlling relaxation dynamics of supercooled liquids---the metastable state of glass-forming liquids---remains a central challenge in condensed matter physics.
The challenge stems from the dramatic non-Arrhenius increase in their relaxation times that can span picoseconds to millions of years upon decreasing temperature \cite{angell2000relaxation}. 
Due to this slowdown, supercooled liquids fall out of equilibrium at the glass transition temperature $T_\mathrm{g}$, limiting access to equilibrium states at ultralow temperatures ($T < T_\mathrm{g}$). 
This inaccessibility prevents investigation of fundamental thermodynamic questions, such as on the existence of the Kauzmann temperature---a transition where non-ergodic glass would emerge as a true equilibrium phase.
When $T < T_\mathrm{g}$, glassy systems display physical aging, where static material properties drift towards the true equilibrium \cite{berthier2011theoretical,biroli2013perspective} on timescales that can even be astronomical REF.
While aging studies have revealed principles such as time reparameterization softness \cite{chamon2002separation,ghimenti2024cleveralgorithmsglassesdo, bohmer2024time}, material-time translational invariance \cite{chamon2002separation,ghimenti2024cleveralgorithmsglassesdo,bohmer2024time}, and violations of the fluctuation-dissipation theorem \cite{chamon2002separation, BarratKob1999, di2000off, grigera1999observation, PhysRevLett.93.160603, wang2006dynamic}, these insights have not resolved the challenge of preparing equilibrium states themselves.

Conventional protocols to prepare supercooled liquids require cooling rates that must be sufficiently slow to maintain adiabatic equilibration, yet fast enough to suppress crystallization---competing requirements that impose severe limitations. %limits.
Only recently have experimental advances, such as physical vapor deposition \cite{ediger2017perspective,rodriguez2022ultrastable}, and computational methods, including swap Monte Carlo algorithms \cite{ninarello2017models,scalliet2022thirty,jung2025numericalinvestigationequilibriumkauzmann}, enabled the preparation of equilibrium states at previously inaccessible temperatures.
However, these methods still fall short of reaching the Kauzmann temperature, rendering fundamental thermodynamics around the glass transition unexplored.


In this Article, we propose a new route to induce controlled aging and cooling in supercooled liquids via strong light–matter coupling. 
We place an atomistic model of a supercooled liquid inside a Fabry–Pérot cavity (\cref{fig:fig1}a), an optical resonator whose mirror spacing $L_\mathrm{c} \sim O(\mu \mathrm{m})$ tunes the resonance frequency $\omega_\mathrm{c}$ and coupling strength $\lambda$. 
The cavity mode hybridizes with molecular vibrations to form polaritons, whose spectroscopic signatures confirm strong light-matter coupling (\cref{fig:fig1}b) \cite{vahala2003optical,aspelmeyer2014cavity}.
Drawing upon theoretical methods from polaritonic chemistry \cite{ebbesen2023introduction,Dunkelberger2022}, we show that the cavity acts as a selective, tunable energy source. 
Unlike a thermal bath that excites all molecular motions, cavity coupling injects energy primarily into dipole-active vibrations.
The sudden increase in intramolecular vibrational energy creates an imbalance that the system restores by driving ro-translational degrees of freedom to deeper potential energy minima (\cref{fig:fig1}c), slowing structural relaxation as the system ages back to equilibrium.
We call this \textit{non-thermal aging}, since velocities remain thermalized at $T$ despite the increased slowdown in structural relaxation.
% (\cref{fig:fig1}d). 
%Thus, the system retains a memory of the protocol on timescales longer than its equilibrium relaxation time.

Nonthermal aging also reveals a fundamental correspondence: nonequilibrium states induced by the cavity at temperature $T$ map to equilibrium states at lower temperatures. 
The correspondence arises because the cavity reparameterizes time evolution through a material time $\xi(t)$, under which nonequilibrium time-correlation functions collapse onto their equilibrium counterparts. 
This collapse implies that aging is a passage through equilibrium states.
To make quantitative use of this mapping, we calculate the structural fictive temperature $T_\mathrm{F}(t)$---the equivalent thermodynamic temperature at which an equilibrium system would exhibit the same intramolecular potential energy as the one measured at time $t$. 
From $T_\mathrm{F}(t)$ and equilibrium relaxation times, we can then use the Tool-Naranaswamy model to quantitatively predict nonthermal aging dynamics at all coupling strengths.

Equipped with this new understanding of cavity aging dynamics, we devise a protocol for preparing equilibrium supercooled liquids, herein referred to as \textit{cavity configurational feedback (C2F) cooling} (\cref{fig:fig1}d). 
Recognizing that realistic cavities exhibit dissipation, we apply time-periodic coupling pulses---a form of Floquet engineering---to systematically quench the structural subsystem despite transient cavity dissipation
Each pulse depresses $T_\mathrm{F}$ to ultra-low temperatures while feedback control adjusts the bath temperature $T$ to track $T_\mathrm{F}(t)$.
Through repeated cycles, we reach a steady state where the periodic forcing ceases. 
At this point, the system remains structurally equilibrated at $T \approx T_\mathrm{g}$—a feat that would require a million years of conventional molecular simulation. 
This achieves our goal of preparing ultralow-temperature equilibrium supercooled liquids via strong light-matter coupling.






\section*{Theoretical Setup}

\paragraph*{The light-matter system.} 
Our system consists of a supercooled liquid confined in a Fabry-Pérot cavity (\cref{fig:fig1}a). 
To target intramolecular vibrations, we set the cavity frequency $\omega_\mathrm{c}$ in the infrared (IR) range, yielding a microcavity geometry. 
The liquid occupies a thin region at the center of the cavity, allowing us to neglect the cavity mode spatial structure. 
Furthermore, we retain only the lowest mode since higher-order modes operate at high frequencies that decouple from any other molecular processes. %the molecular dynamics of interest.
With this setup, the light-matter system consists of $N$ atoms, each with momentum and position $(\bm{P}_i,\bm{R}_i)$ as well as mass $M_i$, and a single cavity mode whose EM field vector along the polarization directions $\{\bm{e}_\alpha\}$ can be assigned to a set of canonical coordinates $(p_\alpha, q_\alpha)$. 

The cavity mode responds to fluctuations in the total molecular dipole moment $\bm{\mu} = \sum_{i=1}^N z_i e \bm{R}_i$, where $z_i$ and $\bm{R}_i$ are the partial charge and position of atom $i$. 
This interaction is captured by the following Hamiltonian:
\begin{equation}
H = H_\mathrm{M}+H_\mathrm{EM}
\end{equation}
where $H_\mathrm{M}$ denotes the molecular Hamiltonian, %which consists of its kinetic and potential energy,
\begin{equation}
H_\mathrm{M}(\{ \bm P_i,\bm R_i \}) = \sum_{i=1}^{N} \frac{|\bm P_i|^2}{2 M_i} +V(\{ \bm R_i\}) \,,
\end{equation}
where $V(\{ \bm R_i \})$ is the molecular potential energy, and $H_\mathrm{EM}$ is the EM field Hamiltonian, which takes the form of a simple harmonic oscillator (SHO)
\begin{equation}
H_\mathrm{EM} (\{ p_\alpha, q_\alpha \}) = \sum_{\alpha=1,2} \left[ \frac{p_\alpha^2}{2}+\frac{1}{2} \omega_\mathrm{c}^2\left(q_\alpha +\frac{ \lambda \mu_\alpha}{\omega_\mathrm{c}} \right)^2 \right], \label{eq:Hem}
\end{equation}
where $\lambda$ is the coupling constant and $\mu_\alpha = \bm \mu \cdot \bm e_\alpha$ is the dipole moment projected on the polarization vectors.
From \cref{eq:Hem}, we see that the interaction corresponds to the displacement of the equilibrium position. % by $\bm q_0 = -\lambda\bm \mu/\omega_\mathrm{c}^2$.

The supercooled liquid follows the Kob-Andersen (KA) model, a glassy two-component Lennard-Jones fluid, extended here to include dipole-active intramolecular vibrations.  
The new molecular KA (mKA) model is a 4:1 mixture of diatomic molecules, $\mathrm{A}_2$ and $\mathrm{B}_2$, with opposing partial charges $\pm 0.3e$ to create permanent dipoles and masses matching an oxygen and nitrogen molecule, respectively.
The intramolecular energy $V_\mathrm{b}$ consists of harmonic terms,
\begin{equation}
V_\mathrm{b} = \sum_{(i,j) \in \mathcal{B}} \frac{1}{2} k_{ij} \left(R_{ij}-R^0_{ij}\right)^2
\end{equation}
where $\mathcal{B}$ is the set of bonds, $R_{ij}=|\bm{R}_i-\bm{R}_j|$ is the distance between atoms $i$ and $j$, $R^0_{ij}$ is the equilibrium bond length. 
The force constant $k_{ij}$ is set so that each molecule has a harmonic vibrational stretch at 1560~cm$^{-1}$ ($\mathrm{A}_2$) or 2433~cm$^{-1}$ ($\mathrm{B}_2$), mimicking  oxygen and nitrogen molecule vibrational frequencies, respectively. 
Intermolecular interactions consist of Coulombic and Lennard-Jones potentials:
\begin{gather}
V_\mathrm{C} = \sum_{(i,j) \in \mathcal{N}} \frac{z_i z_j e^2}{4 \pi \epsilon_0 R_{ij}}, \\
V_\mathrm{LJ} = \sum_{(i,j) \in \mathcal{N}} 4 \epsilon_{ij}\left[\left(\frac{\sigma_{ij}}{R_{ij}}\right)^{12}-\left(\frac{\sigma_{ij}}{R_{ij}}\right)^6\right]
\end{gather}
where $\mathcal{N}$ is the set of all atom pairs excluding bonded pairs, and thus the total potential energy $V= V_\mathrm{b} + V_\mathrm{LJ} +V_\mathrm{C}$.
LJ parameters for the $\mathrm{A}_2$ molecule ($\lambda_\mathrm{AA}$ and $\sigma_\mathrm{AA}$) are tuned to those of an O$_2$ molecule, and the remaining parameters, including cut-off radius,  are scaled to preserve the original KA model characteristics. 
%These ingredients yield a glassy system with a simple vibrational structure, making it ideal for exploring strong light-matter coupling in a microcavity.

\paragraph*{Equations of motion (EOM).} 
Dynamics follows classical Hamilton's equations with additional forces from the external baths.
The cavity mode exchange energy with a Langevin bath on timescales $\tau_\mathrm{c}$ so that its EOM becomes
\begin{gather}
\dot{q}_{\alpha} = p_\alpha \\
\dot{p}_\alpha = -\omega_\mathrm{c}^2 \left(q_\alpha + \frac{\lambda \mu_\alpha}{\omega_c^2} \right) -\gamma_\mathrm{c} p_\alpha +\sqrt{2 \gamma_\mathrm{c} k_\mathrm{B} T } \eta_\alpha(t)
\end{gather}
where $\eta_\alpha(t)$ is white noise, $\gamma_\mathrm{c} = 1/(2\tau_\mathrm{c})$. 
% Outside the cavity, the supercooled liquid is also thermalized by an external bath at timescales $\tau_\mathrm{b}$ and temperature $T$. 
% Unlike the cavity bath, the molecular bath should not disturb the microscopic dynamics inherent to the glassy system, prompting the use of as a Bussi-Parrinello bath model, which exchanges energy with the molecular system like a Langevin bath while minimally perturbing  dynamics. 
% The EOM for the molecular system then becomes 
Outside the cavity, the supercooled liquid is thermalized by an external Bussi-Parrinello bath at temperature $T$ and timescale $\tau_\mathrm{b}$, which exchanges energy like a Langevin thermostat while minimally perturbing the microscopic glassy dynamics. The EOM for the molecular system is 
\begin{gather}
\dot{\bm{R}}_i = \frac{\bm{P}_i}{M_i} \\
\dot{\bm{P}}_i = -\nabla_{\bm{R}_i} V -\gamma(K) \bm P_i + \sqrt{2 D(K)} \bm P_i \eta_i (t)
\end{gather}
where $K$ is the total kinetic energy, and the generalized friction and diffusion coefficients are $\gamma(K) = \frac{1}{2\tau_\mathrm{b}} \left(1-\left(1-\frac{1}{3 N}\right) \frac{\bar{K}}{K}\right)$ and $ D (K) = \bar{K} / (6 N \tau _\mathrm{b} K)$, respectively, where $\bar{K} = \frac{3}{2}Nk_\mathrm{B} T$. This fully classical treatment of our system enables molecular dynamics (MD) simulations to many-molecule response and relaxation of our system over long timescales. 
Following Ref.~X, we call this approach \textit{cavity molecular dynamics} (cavMD). Details of the Hamiltonian, equations of motion, and propagation scheme are given in \cref{app:theory}.


\begin{figure*}[t]
    \centering
    \includegraphics[width=0.9\linewidth]{Figure2.pdf}
    \caption{Cavity-induced slowdown of structural relaxation and nonthermal aging dynamics. 
    (a) Top row: Infrared absorption spectra $n(\omega)\alpha(\omega)$ showing polariton splitting at different coupling strengths $\lambda$. Gray dashed line indicates cavity frequency $\omega_\mathrm{c} = 1560$ cm$^{-1}$. Bottom row: Normalized intermediate scattering functions $\phi_{\bm{k}}(t;t_\mathrm{w})$ at multiple waiting times $t_\mathrm{w}$ (color scale from purple to yellow: $0$ to $2400$ ps) for each coupling strength. Stronger coupling induces greater initial slowdown, with dynamics gradually recovering toward equilibrium at longer $t_\mathrm{w}$.
    (b) Normalized structural relaxation times $\tilde{\tau}_\mathrm{s} = \tau_\mathrm{s}/\tau_{s,\lambda=0}$ versus coupling strength $\lambda$ at different waiting times (color-coded by $t_\mathrm{w}$), where $\tau_{s,\lambda=0}$ is the equilibrium relaxation time at $\lambda = 0$. The slowdown effect increases with $\lambda$, saturating at $\tilde{\tau}_\mathrm{s} \approx 4.75$ for $\lambda = 0.141$ a.u. at $t_\mathrm{w} = 0$, and gradually decreases toward equilibrium (yellow-green curves) at large $t_\mathrm{w}$.
    (c) Time evolution of $\tilde{\tau}_\mathrm{s}$ over waiting times for different coupling strengths. Memory effects persist up to $\sim 2.5$ ns ($\sim 50\times$ the equilibrium relaxation time), with stronger coupling showing larger initial slowdown and longer relaxation back to equilibrium.}
    \label{fig:figure2}
\end{figure*}

\begin{figure*}[t]
    \centering
    \includegraphics[width=\linewidth]{Figure3.pdf}
    \caption{
    (a) Hierarchy of characteristic timescales: Rabi oscillations ($\Omega_R^{-1}$, femtoseconds) drive energy exchange between cavity and molecular vibrations; molecular and cavity baths ($\tau_\mathrm{b}$, $\tau_\mathrm{c}$, picoseconds) dissipate energy to keep overall equilibrium at all times; structural relaxation ($\tau_\mathrm{s}$, nanoseconds) governs glassy dynamics, with $\tau_\mathrm{s} \gg \tau_\mathrm{b} \gg \Omega_R^{-1}$.
    (b) Time evolution of potential energy components following cavity activation at $t_0 = 200$ ps with $\lambda = 0.141$ a.u. Vibrational energy (blue) spikes initially due to cavity pumping, while Lennard-Jones and Coulombic energy (red) drops instantaneously to maintain near-equilibrium total potential energy (green). Dashed line indicates equilibrium reference.
    (c) Schematic illustration of two types of fictive temperatures, the vibrational ($T_\mathrm{v}$) and structural ($T_\mathrm{s}$) fictive temperature.
    (d) Fictive temperature evolution for $\lambda = 0.141$ a.u. Vibrational temperature $T_\mathrm{v}$ (blue) reaches $\sim 800$ K while overall temperature (green) remains near 100 K. Structural fictive temperature $T_\mathrm{F}$ (red) drops to $\sim 70$ K to compensate, maintaining thermal equilibrium. Inset: detailed view of initial 600 ps showing rapid onset and gradual relaxation back to equilibrium.
    (e) Initial structural fictive temperature $T_\mathrm{F}^0$ immediately after cavity activation versus coupling strength $\lambda$. Self-quenching effect increases with $\lambda$.}
    \label{fig:fig3}
\end{figure*}


\paragraph*{Observables.} 
To probe dynamics, we can calculate observables such as dipole moments with the time-correlation function being
\begin{equation}
C_\mu (t) = \sum_{\alpha=1,2} \left\langle \mu_\alpha(t) \mu_\alpha(0) \right\rangle \,,
\end{equation}
where $\langle \ldots \rangle$ denotes ensemble averaging. Its Fourier transform yields the IR absorption coefficient $\alpha(\omega)$, which is sensitive to polariton formation,
\begin{equation}
n(\omega) \alpha(\omega) = \frac{\beta \omega^2}{4 \epsilon_0 V c} \int_{-\infty}^{+\infty} \!\!\!\! dt \, e^{-i \omega t} C_\mu(t)
\end{equation}
where $n(\omega)$ is the refractive index and $\beta = 1/k_\mathrm{B} T$. At equilibrium ($T = 100$ K, \cref{fig:fig1}b), tuning $\omega_\mathrm{c} = \omega_\mathrm{AA}$ splits the vibrational peak into upper and lower polariton branches separated by the Rabi splitting $\Omega_\mathrm{R}$, which characterizes the frequency of energy exchange within our light-matter system and increases linearly with $\lambda$.

For structural relaxation at long timescales, we analyze the density mode $\rho_{\bm{k}}(t) = \sum_{i=1}^N e^{i \bm{k} \cdot \bm{R}_i(t)}$ of wavevector $\bm k$ and its time correlation function
\begin{equation}
F_{\bm{k}}(t) = \langle \rho_{\bm{k}}(t)\rho^*_{\bm k}(0) \rangle \,.
\end{equation}
Also known as the intermediate scattering function (ISF),  $F_{\bm k}(t)$ captures collective molecular rearrangement and is measurable via neutron or X-ray scattering. Its value at $t=0$ defines the static structure factor $S_{\bm k}= \langle \rho_{\bm{k}}(0)\rho^*_{\bm k}(0) \rangle$, which we use to normalize the ISF as 
\begin{equation}
\phi_{\bm k}(t) := F_{\bm{k}}(t)/S_{\bm k}\,,
\end{equation}
and the relaxation time $\tau_\mathrm{s}$ is defined as $\phi_{\bm k}(\tau_\mathrm{s}) = 0.1$.
%At equilibrium, structural relaxation inside the microcavity is indistinguishable from outside at all temperatures (Fig.~SREF), consistent with previous cavMD simulations.

\paragraph*{Cavity aging protocol.} 
In conventional aging, a supercooled liquid equilibrated at $T_i$ is quenched to $T_f < T_i$, and properties slowly relax toward equilibrium. In this nonequilibrium regime, correlation functions depend on both elapsed time $t$ and the elapsed time after quench, referred to as the waiting time $t_\mathrm{w}$. Here, the ISF becomes
\begin{equation}
F_{\bm{k}}(t; t_\mathrm{w}) = \langle \rho_{\bm{k}}(t+t_\mathrm{w}) \rho^*_{\bm k}(t_\mathrm{w}) \rangle \,.
\end{equation}
The static structure factor becomes time-dependent, $S_{\bm k}(t_\mathrm{w}) = \langle \rho_{\bm k}(t_\mathrm{w}) \rho^*_{\bm k}(t_\mathrm{w}) \rangle$, yielding a new normalized ISF:
\begin{equation}
\phi_{\bm k}(t; t_\mathrm{w}) = F_{\bm{k}}(t; t_\mathrm{w})/S_{\bm k}(t_\mathrm{w}) \,,
\end{equation}
and thus the relaxation time $\tau_\mathrm{s}(t_\mathrm{w})$ also acquires waiting-time dependence.
Inspired by these aging experiments, we introduce an analogous cavity aging experiment where, instead of changing temperature, we activate light-matter coupling at time $t_0$:
\begin{equation}
\lambda(t) = \lambda \Theta(t-t_0) \,. \label{eq:varepsilon}
\end{equation}
Similar to a thermal aging experiment, material properties should also age toward a new equilibrium, but unlike thermal quenches, the bath temperature $T$  remains constant throughout.

% \begin{equation}
% F_{\bm{k}}(t; t_\mathrm{w}) = \langle \rho_{\bm{k}}(t+t_\mathrm{w}) \rho^*_{\bm k}(t_\mathrm{w}) \rangle ,
% \end{equation}

% with ensemble averages taken under the quench protocol.

% Inspired by this, we define a cavity aging experiment (\cref{fig:fig1}d), where instead of a thermal quench we switch on light–matter coupling at $t_0$:

% \begin{equation}
% \lambda(t) = \lambda \Theta(t-t_0) , \label{eq:varepsilon}
% \end{equation}
%so the system ages toward a new equilibrium under constant bath temperature $T$.

% In conventional aging experiments, a supercooled liquid equilibrated at initial temperature $T_i$ is rapidly quenched to final temperature $T_f < T_i$. 
% Material properties such as specific heat, enthalpy, and viscosity evolve then slowly toward equilibrium at $T_f$. 
% In terms of our observables, the aging evolution manifests as violations of time-translational invariance in time-correlation functions since the system is not in a stationary state.
% Thus, during an aging experiment, we must track the time origin when we perform our measurements, which is often referred to as the waiting time $t_\mathrm{w}$.
% Taking the ISF as an example, the dependence of waiting time can be measured through 

% In conventional aging, a supercooled liquid equilibrated at $T_i$ is quenched to $T_f < T_i$, and properties like specific heat, enthalpy, and viscosity slowly relax toward equilibrium at $T_f$. In this nonequilibrium regime, correlation functions cannot be written as functions of time differences alone. Instead, they also depend on the time elapsed since the quench, called the waiting time $t_\mathrm{w}$. For example, the ISF depends on $t_\mathrm{w}$ via
% \begin{equation}
% F_{\bm{k}}(t; t_\mathrm{w}) = \langle \rho_{\bm{k}}(t+t_\mathrm{w}) \rho^*_{\bm k}(t_\mathrm{w}) \rangle
% \end{equation}
% where $\langle \ldots \rangle$ now denotes nonequilibrium ensemble averaging under the quench or other nonequilibrium protocol.
% The static structure factor now acquire time dependence with $S_{\bm k}  (t_\mathrm{w}) = \langle \rho_{\bm k}(t_\mathrm{w}) \rho^*_{\bm k}(t_\mathrm{w}) \rangle$, and thus the normalized ISF becomes:
% \begin{equation}
% \phi_{\bm k}(t; t_\mathrm{w}) = F_{\bm{k}}(t; t_\mathrm{w})/S_{\bm k}(t_\mathrm{w}) 
% \end{equation}
% , we introduce an analogous cavity aging experiment (\cref{fig:fig1}d) where, instead of changing temperature, we activate light-matter coupling at time $t_0$, i.e., 
% \begin{equation}
% \lambda(t) = \lambda \Theta(t-t_0) \,, \label{eq:varepsilon}
% \end{equation}
% where $\Theta(t)$ is the Heaviside function. \Cref{eq:varepsilon} models the effect when we first expose the supercooled liquid to a cavity mode. 
% Similar to a thermal aging experiment, material properties should also age toward a new equilibrium, but unlike thermal quenches, the bath temperature remains constant at $T$ throughout. 

\paragraph*{Computational details.} 
In what follows, we present results of our cavity aging experiments, and discuss their broad implications in manipulating relaxation of supercooled liquids. All results are reported in atomic units unless explicitly noted. 
We simulate $N = 250$ molecules at density $\rho = 0.03125$ in a periodic box of side length $L = 40$. 
We choose $T=100$ K where the system is modestly supercooled with equilibrium relaxation time $\tau_\mathrm{s} \sim 100$ ps---roughly an order of magnitude longer than at onset temperature $T_\mathrm{o} = 134.1$ K where non-Arrhenius behavior begins.
We explore coupling strengths $\lambda = (0.042, 0.070,0.098, 0.141)$ a.u. with cavity frequency $\omega_\mathrm{c} = 1560$ cm$^{-1}$ matching the $\mathrm{A}_2$ stretch.
Both molecular and cavity bath relaxation times are $\tau_\mathrm{m} = \tau_\mathrm{c} = 1.0$ ps, slower than vibrational timescales.
In our cavity aging protocol, we activate coupling at $t_0 = 200$ ps, with measurements at waiting time intervals $\Delta t_\mathrm{w} =200$ ps up to $t_\mathrm{w} = 2.6$ ns.
Ensemble averaging employs $N_\mathrm{traj} = 500$ independent trajectories.
Simulations use \texttt{cavHOOMD-blue} REF, our GPU-accelerated extension of \texttt{HOOMD-blue} incorporating light-matter interactions (details in Appendix~REF).

\section*{Nonthermal Aging}

\paragraph*{Results} \Cref{fig:fig1}e shows a comparison of equilibrium and nonequilibrium ISFs at three waiting times for coupling $\lambda = 0.098$ a.u., revealing the sudde slowdown of structural relaxation during nonthermal aging. 
The full extent of how the coupling strength reshapes the ISF can be found in \cref{fig:figure2}a, where we see that stronger coupling continuously extends relaxation to longer timescales. 
From \cref{fig:figure2}b, we also see that this effect saturates, with $\tau(t_\mathrm{w}=0)$ reaching maximum slowdown of $4.75\times$ at $\lambda = 0.141$ a.u.
The increased slowdown persists over extended waiting times (\cref{fig:figure2}c), with memory effects lasting up to 2.5 ns—roughly $50\times$ the equilibrium relaxation time.
These results are noteworthy given that energy transfer from the cavity occurs on timescales far shorter than structural relaxation, yet can leave lasting memory effects on glassy dynamics.

To understand non-thermal aging, \cref{fig:fig3}a illustrates the hierarchy of timescales governing energy transfer: Rabi oscillations between molecular vibrations and cavity modes occur on femtosecond scales ($\Omega_R^{-1}$), energy dissipation to thermal baths on picosecond scales ($\tau_\mathrm{b}$), and structural relaxation on nanosecond scales ($\tau_\mathrm{s}$), such that $\tau_\mathrm{s} \gg \tau_\mathrm{b} \gg \Omega_R^{-1}$.
At $t = t_0$, the cavity rapidly transfers energy to molecular vibrations on the Rabi timescale, creating a spike in vibrational energy (\cref{fig:fig3}b).
Thermal dissipation maintains the total potential energy at near equilibrium, forcing translational and rotational degrees of freedom to compensate by accessing lower potential energy configurations.
This self-quenching is evident in \cref{fig:fig3}b, where the total Lennard-Jones and Coulombic energies, which are sensitive to ro-translational motion, drop nearly instantaneously on subpicosecond timescales.
The system then gradually ages back to equilibrium on the structural relaxation timescale.

% To understand how non-thermal aging works, \cref{fig:fig3} illustrates the key processes that govern energy transfer in the system: Rabi oscillations between molecular vibrations and cavity modes occur on the fastest femtosecond timescale ($\Omega_R^{-1}$), energy dissipation to thermal baths operates on intermediate picosecond scales ($\tau_\mathrm{b}$), while structural relaxation of the supercooled liquid extends to the longest nanosecond timescales ($\tau_\mathrm{s}$), such that $\tau_\mathrm{s} \gg \tau_\mathrm{b} \gg \Omega_R^{-1}$.
% This hierarchy of timescale separations results in the following process: 
% at $t = t_0$, the microcavity rapidly transfers energy to the molecular vibrations on the Rabi timescale.
% As \cref{fig:fig3}b demonstrates, this initial energy transfer creates a spike in vibrational mode energy.
% Subsequently, on the bath timescale, the molecular bath dissipates excess energy to maintain the total potential energy close to equilibrium.
% However, since the vibrational vibrational modes still retain excess energy while the overall energy remains close to equilibrium, the translational and rotational degrees of freedom must compensate by accessing configurations with lower potential energy.
% This compensation is evident in \cref{fig:fig3}b, where the Lennard-Jones and Coulombic potential energies drop simultaneously with the vibrational energy increase.
% This self-quenching mechanism drives the system into deeper potential energy minima, making the system more glassy before gradual aging on the structural relaxation timescale returns the structure to equilibrium.

\paragraph*{Fictive temperatures.} We can gain a qualitative understanding of nonthermal aging by translating potential energies into effective temperatures.
To this end, we introduce the concept of a fictive temperature, defined as the equilibrium temperature that would yield the instantaneous potential energy measured at time $t$.
Moreover, timescale separation allows us to treat harmonic vibrations and structural relaxation as separate subsystems with distinct fictive temperatures.
Let the equilibrium vibrational energy be $\varphi(T):= \langle V_\mathrm{b} \rangle $. Equipartition gives $\varphi(T) = \frac{1}{4} N k_\mathrm{B} T$ for $T \to 0$ and assuming $\varphi(T) \to \varphi_\infty$ at high $T$, we can write the following Pad\'e approximant:
\begin{equation}
    \varphi(T) = \frac{c_\mathrm{v} T}{1+(c_\mathrm{v} /\varphi_\infty)T  } \,, \quad c_\mathrm{v}=\frac{1}{4} N k_\mathrm{B}
\end{equation}
and the vibrational fictive temperature can be written as
\begin{equation}
T_\mathrm{v}(t) := \varphi^{-1}[V_\mathrm{b}(t)] \,. \label{eq:Tv}
\end{equation}
%\Cref{eq:Tv} is useful as a diagnostic; when $T_\mathrm{v}(t) > T$, the system exhibits large-amplitude vibrations beyond equilibrium fluctuations.
To track structural degrees of freedom, we can use the equilibrium intramolecular potential energy  $\vartheta(T) :=\langle V_\mathrm{LJ}+V_\mathrm{C}\rangle $, whose low-temperature limit follows the Rosenfeld-Tarazona scaling, $\vartheta(T) = \vartheta_0+\alpha_0 T^{3/5}$. 
Assuming $\vartheta(T) \to \vartheta_\infty$ for high-$T$, we can also construct the following Pad\'e approximant: % accurate across all temperatures:
\begin{equation}
\vartheta(T) = \vartheta_0+\frac{\alpha_0 T^{3/5}}{1+\frac{\alpha_0}{\vartheta_\infty-\vartheta_0}T^{3/5}} \,. \label{eq:varthetaT}
\end{equation}
Using \cref{eq:varthetaT}, we now define the structural fictive temperature as
\begin{equation}
T_\mathrm{F}(t) := \vartheta^{-1}[V_\mathrm{LJ}(t)+V_\mathrm{C}(t)] \,.
\end{equation}
Since Lennard-Jones and Coulombic interactions govern glassy dynamics, $T_\mathrm{F}$ also indicates the system depth in the PEL in terms of a temperature scale (\cref{fig:fig3}c).

Applying these definitions to the measured potential energies, we obtain \cref{fig:fig3}d, which shows the fictive temperatures of different subsystems during cavity-induced aging, alongside the kinetic temperature $T_\mathrm{k}$. 
During energy pumping, $T_\mathrm{v}$ reaches $\sim$800~K while the $T_\mathrm{k}$ remains near 100~K.
At the highest coupling strength, the structural degrees of freedom quenches down, with $T_\mathrm{F}$ approaching $\sim$70~K, to maintain constant kinetic temperature, before both subsystems eventually equilibrate over nanosecond timescales. 
The extent of in which a cavity reduces $T_\mathrm{F}$ can be found in \Cref{fig:fig3}e, where we observe that the self-quenching effect increases with $\lambda$ but saturates at strong coupling, reaching a maximum depression of $\Delta T_\mathrm{F} \approx 30$ K at $\lambda = 0.141$ a.u. 

\begin{figure}[t]
    \centering
    \includegraphics[width=\linewidth,height=0.65\linewidth]{example-image-a}
    \caption{A short plot on polaritons and being off-resonant. (a) It will have the plot as a function of frequency and then (b) It will have the plot as a function of waiting time. A few more comments.}
    \label{fig:placeholder}
\end{figure}

\begin{figure*}[t]
    \centering
    \includegraphics[width=0.875\linewidth]{Figure4.pdf}
    \caption{Validating time reparameterization softness and predicting aging dynamics. 
    (a) Material time $\xi(t)$ reconstructed from ISF measurements (circles) at different coupling strengths $\lambda$ (color-coded). Stronger coupling leads to slower accumulation of material time.
    (b) Collapse of all normalized ISFs onto a universal master curve following stretched exponential $\Phi_{\bm{k}}(\xi) = e^{-\xi^\beta}$ with $\beta = 0.55$ (black line), independent of coupling strength and waiting time. This collapse confirms material-time translational invariance and demonstrates that all nonequilibrium states, including equilibrium ($\lambda = 0$), share the same functional form.
    (c) The Tool--Narayanaswamy (TN) model predictions  for material time, $\xi_\mathrm{TN}(t)$ (dashed lines) computed from equilibrium relaxation times and measured fictive temperatures.
    (d) TN prediction for normalized relaxation times $\tilde{\tau}_\mathrm{s}$ versus coupling strength $\lambda$ at different waiting times $t_\mathrm{w}$ (color scale). TN predictions (lines with diamonds) successfully reproduce the trends observed in \cref{fig:figure2}b, confirming that equilibrium properties and fictive temperature fully determine aging dynamics.
    (e) TN prediction for time evolution of $\tilde{\tau}_\mathrm{s}$ over waiting times for different coupling strengths, accurately capturing the relaxation back to equilibrium.}
    \label{fig:fig4}
\end{figure*}

\paragraph*{The role of polaritons} [Wait for Arianna's' results,
Hopefully this can be a short paragraph on whether it's more like a finite-q effect, poolartionic, etc. 
\lipsum[4]]

\begin{figure*}[t]
    \centering
    \includegraphics[width=\linewidth,height=0.4\linewidth]{Figure6.pdf}
    \caption{Summary of cavity configurational feedback (\texttt{C2F}) cooling of supercooled liquids.}
    \label{fig:fig6}
\end{figure*}


\section*{Time Reparameterization via Material Time}

\paragraph*{Time reparameterization softness.} 
The fictive temperatures we measure provide a quantitative framework for predicting aging dynamics. 
To establish this, we verify whether cavity-induced nonthermal aging obeys time-reparameterization softness (TRS), a universal feature of glassy dynamics.
TRS states that out-of-equilibrium correlation functions evolve by changing their pace rather than their functional form. 
This concept is formalized through material time $\xi(t)$, which ticks at a rate, possibly nonlinear in $t$, as the system ages.
Under TRS, correlation functions obey material-time translational invariance (MTTI), which for the ISF implies
\begin{equation}
\phi_{\bm{k}}(t; t_\mathrm{w}) = \Phi_{\bm k}(\xi(t+t_\mathrm{w})-\xi(t_\mathrm{w})) \label{eq:MTTI}
\end{equation}
where $\Phi_{\bm k}$ is a master function independent of $t_1$ and $t_2$. 
At equilibrium, $\xi(t) = t/\tau_\mathrm{s}$ recovers usual time-translational invariance.
Out of equilibrium, we can use \cref{eq:MTTI} to reconstruct $\xi(t)$ from measurements at multiple waiting times; see Appendix~REF details the statistical procedure.
\Cref{fig:fig3}d shows the reconstructed $\xi(t)$, revealing that stronger coupling slows down material time accumulation. 
\Cref{fig:fig3}e demonstrates collapse of all relaxation functions onto a universal stretched exponential
\begin{equation}
\Phi_{\bm k}(\xi) = e^{-\xi^\beta}
\end{equation}
with $\beta = 0.55$, independent of coupling strength, confirming that cavity aging is controlled by the system's inherent glassiness. 

\paragraph*{Predicting aging dynamics.} 
Since \cref{fig:fig3}e includes the zero-coupling equilibrium case, all nonequilibrium ISFs collapse onto a single equilibrium master curve. 
This collapse reveals a fundamental correspondence: nonequilibrium states from aging maps exactly onto a sequence of equilibrium states at progressively lower temperatures.
We test this equilibrium correspondence by predicting aging dynamics entirely from equilibrium properties using the Tool--Narayanaswamy (TN) phenomenological model. 
In the TN model, material time evolves at a rate set by the equilibrium relaxation time at the instantaneous fictive temperature:
\begin{equation}
    \frac{d \xi_\mathrm{TN}(t)}{d t}  = \frac{1}{\tau_\mathrm{s}[T_\mathrm{F}(t)]}
\end{equation}
Combined with \cref{eq:MTTI}, this enables complete reconstruction of aging trajectories from equilbrium properties and fictive temperature measurements.

\Cref{fig:fig4}a demonstrates quantitative agreement between TN-predicted $\xi_\mathrm{TN}(t)$ and direct measurements across all coupling strengths.
The inset shows the aging rate $d\xi/dt$ starting lower due to initial fictive temperature depression and gradually recovering as the system ages back toward equilibrium.
\Cref{fig:fig4}b,c further validate the mapping: TN predictions successfully reproduce both coupling-strength dependence and waiting-time evolution of normalized relaxation times.
These results confirm that cavity aging explores a thermally accessible trajectory through equilibrium phase space, albeit one reached through nonthermal means.

\section*{Cavity Configurational Feedback (\texttt{C2F}) Cooling}

The fundamental difficulty in preparing equilibrium supercooled liquids stem from competing timescales. As temperature decreases, relaxation times increase exponentially, eventually exceeding experimental timescales at the glass transition temperature $T_\mathrm{g}$. Below $T_\mathrm{g}$, systems fall out of equilibrium: while the kinetic temperature $T_\mathrm{k}$ rapidly tracks the bath temperature through fast energy dissipation (picosecond timescales), the fictive temperature $T_\mathrm{F}$ lags far behind due to slow structural rearrangements that can span astronomical timescales. Conventional thermal protocols cannot cool the structural subsystem fast enough to maintain equilibrium.

Our aging experiments reveal a solution: cavity coupling actively drives supercooled liquids into deeper potential energy minima, accessing equilibrium configurations at lower fictive temperatures. Rather than waiting for thermal fluctuations, we exploit periodic cavity pumping—a form of Floquet engineering—to systematically quench intermolecular degrees of freedom. 
To this end, we develop \textit{cavity configurational feedback (\texttt{C2F}) cooling}---a protocol that drives fictive temperature toward target values while maintaining structural equilibrium through time-periodic energy injection.

\paragraph*{Protocol design.} 
\texttt{C2F} cooling combines two key elements (\cref{fig:fig6}a): structural thermometry measuring $T_\mathrm{F}(t)$ at picosecond resolution and ultrafast feedback control adjusting bath temperature dynamically. 
In simulations, we achieve structural thermometry by utilizing the analytical form of the potential energy directly; experimentally, spectroscopic probes of intermolecular structure could serve this role, analogous to Raman thermometry for vibrations. The feedback control law
\begin{equation}
\frac{dT(t)}{dt} = -\frac{1}{\tau_\mathrm{f}}\left(T(t)-T_\mathrm{F}(t) \right) \,,
\end{equation}
allows the bath to track the fictive temperature at timescale $\tau_\mathrm{f}$. Simultaneously, we apply periodic cavity pulses with exponential decay modeling realistic photon loss:
\begin{equation}
\lambda(t) = \lambda_0 \exp \left[-\frac{t- \lfloor \Omega_\mathrm{p} t\rfloor \Omega^{-1}_\mathrm{p}}{\tau_\mathrm{p}}\right]
\end{equation}
where $\tau_\mathrm{p} = 0.1$ ps and $\Omega_\mathrm{p}^{-1} = 100$ ps. This time-periodic forcing constitutes the Floquet engineering of the supercooled liquid: each pulse transiently injects energy into intramolecular vibrations, creating an imbalance that drives structural degrees of freedom into deeper minima repeatedly. Feedback control then lowers the bath temperature $T$ to match the depressed $T_\mathrm{F}$. Through repeated cycles, the system reaches a non-equilibrium steady state where structural degrees of freedom equilibrate ($\approx T \approx T_\mathrm{F}$).

\paragraph*{Results.} To demonstrate \texttt{C2F} cooling, we compare two protocols starting from equilibrium at $T_\mathrm{initial} = 100$ K. Conventional thermal cooling uses feedback control alone:
\begin{equation}
\frac{dT(t)}{dt} = -\frac{1}{\tau_\mathrm{f}} \left(T(t)-T^\star\right)
\end{equation}
Without cavity assistance, \cref{fig:fig5}c reveals severe lag: while $T_\mathrm{k}$ reaches 5 K, the fictive temperature remains at $T_\mathrm{F} \approx 20$--25 K, indicating the structure has fallen far out of equilibrium.

In contrast, \texttt{C2F} cooling with periodic cavity pulses at $\lambda_0 = 0.141$ a.u. achieves structural equilibration. \Cref{fig:fig5}b shows the time evolution of all three temperatures during \texttt{C2F} cooling. Each cavity pulse creates a transient depression in $T_\mathrm{F}$, which feedback control matches by lowering the bath temperature. Through repeated pulses, $T_\mathrm{F}$ progressively decreases until reaching a steady state where cavity-induced quenching balances structural relaxation. The inset of \cref{fig:fig5}b demonstrates that the final achieved temperature depends systematically on coupling strength, with $\lambda_0 = 0.141$ reaching $T_\mathrm{F,final} \approx 8$ K. Critically, this steady state is a non-equilibrium one: while $T_\mathrm{F} \approx T \approx T_\mathrm{k}$ (structural equilibrium), the vibrational temperature remains elevated at $T_\mathrm{v} \approx 200$ K due to continuous periodic driving—a hallmark of Floquet-engineered systems.

\Cref{fig:fig5}c compares final fictive temperatures against target temperatures for both protocols. The benchmark protocol shows large deviations from the diagonal $T_\mathrm{F} = T_\mathrm{target}$, confirming failure to maintain equilibrium. \texttt{C2F} cooling tracks the diagonal closely for $T_\mathrm{target} > T_\mathrm{g}$, demonstrating successful structural equilibration. Below $T_\mathrm{g} \approx 50$ K, even \texttt{C2F} shows some deviation, but achieves fictive temperatures substantially lower than conventional thermal protocols---accessing structurally equilibrated states at temperatures previously unreachable on computational timescales.


\paragraph*{Post-protocol stability.} After achieving the target fictive temperature, we deactivate periodic forcing and monitor long-term stability. The vibrational modes retain excess energy ($T_\mathrm{v} > T$) accumulated during the Floquet driving phase. This stored energy dissipates through the glassy subsystem, potentially reheating the supercooled liquid and undoing the configurational cooling. 


\Cref{fig:fig5}d addresses this concern by comparing temperature evolution after deactivating cavity pulses under two conditions: (1) constant bath temperature, and (2) continued feedback control following Eq.~REF. With constant bath, we observe modest reheating as vibrational energy slowly dissipates ($\Delta T_\mathrm{F} \approx 2$--3 K over 1 ns). However, maintaining active feedback control successfully stabilizes all three temperatures at the target value, demonstrating that structurally equilibrated low-$T_\mathrm{F}$ states can be preserved indefinitely with appropriate thermal management even after ceasing Floquet driving.


\section*{Discussions}

%\section*{Towards the Kauzmann Transition}

\bibliography{apssamp}

\appendix 

% \section{Deriving the Light-Matter Hamiltonian}

% In this section, we sketch a self-contained derivation of the light-matter Hamiltonian used in this study. To begin, we write the minimal coupling Hamiltonian $\hat{H}$ describing how matter, consisting of negatively and positively charged electrons and nuclei, respectively, interacts with an electromagnetic (EM) field
% \begin{gather}
% \hat{H} = \hat{H}_\mathrm{M}+\hat{H}_\mathrm{F} \label{eq:H}
% \\
% \hat{H}_\mathrm{M} = \sum_{i=1}^{N_\mathrm{n}} \frac{| \hat{\bm P}_i -Z_i \hat{\bm A}(\bm R_i)|^2}{2 M_i} + \sum_{i=1}^{N_e} \frac{| \hat{\bm p}_i -e \hat{\bm A}(\bm r_i)|^2}{2 m_i} \nonumber
% \\
% + \hat{V}_\mathrm{nn}+ \hat{V}_\mathrm{ee}+ \hat{V}_\mathrm{ne} \label{eq:HM}
% \\
% \hat{H}_\mathrm{F} = \frac{\epsilon_0}{2}\int_\Omega d^3 \bm r  \ \big( | \hat{\bm E}(\bm r)|^2+c^2 | \hat{\bm  B}(\bm r)|^2 \big) \label{eq:HF}
% \end{gather}
% where $\hat{H}_\mathrm{F}$ and $\hat{H}_\mathrm{M}$ are the light and matter parts of the Hamiltonian. The first two terms in $\hat{H}_\mathrm{M}$ represent the kinetic energy of the nuclei and electrons, coupled to the vector potential $\bm$ of the EM field, and the rest of the terms represent nucleus–nucleus, electron–electron, and nucleus–electron Coulombic interactions. 

% In writing \cref{eq:H,eq:HM,eq:HF}, we have already made a number of assumptions. First, we have assumed that the EM field is described in the Coulomb gauge, where $\nabla \cdot \mathbf{A} = 0$, and we have subsumed the scalar potential $\phi$ into static Coulombic interactions of charged particles. Thus, the remaining EM fields written in \cref{eq:H,eq:HM,eq:HF} are the purely transverse fields present in the system. 

% To cast \cref{eq:H,eq:HM,eq:HF}, we begin writing the light Hamiltonian in the familiar harmonic-oscillator form. Here, we can work classically first and write the energy density $u$ of the EM field in terms of the vector potential, making use of the fact that $\bm E = -\partial_t \bm A$ and $\bm B = \nabla \times \bm A$:
% \begin{equation}
% H_\mathrm{EM}  = \frac{\epsilon_0}{2} \left(\left|\partial_t \bm A \right|^2
% + c^2 |\nabla \times \bm A|^2 \right)
% \end{equation}
% For a concrete illusration, we can assume 1D Fabry--Perot cavity bounded at $z=\pm L/2$, we can expand the transverse vector potential in terms of normal mode functions
% \begin{equation}
% \mathbf A(\bm r,t) = \sum_{n,\lambda} q_{n,\lambda}(t)\, \bm u_{n,\lambda}(\bm r),
% \
% \bm u_{n,\lambda}(\bm r)= u_n(z)\,\hat{\mathbf e}_\alpha,
% \end{equation}
% where $\hat{\mathbf e}_\alpha\!\perp\!\hat{\mathbf z}$ are fixed polarization vectors and
% \begin{equation}
% u_n(z) = \sqrt{\frac{2}{V}}\,\sin \left(\frac{n\pi}{L}\left(z+\frac{L}{2}\right)\right),
% \quad \omega_n=\frac{n\pi c}{L}.    
% \end{equation}
% where $V=A L$. The mode functions satisfy
% \begin{equation}
% \int_V d^3\bm r\, \bm u_{n,\lambda}\cdot \bm u_{m,\lambda'} = \,\delta_{nm}\delta_{\lambda\lambda'}.    
% \end{equation}
% Substituting this expansion into the Hamiltonian and using orthogonality, each mode decouples:
% \begin{equation}
% H_\mathrm{EM} = \frac12 \sum_{n,\lambda} \Big( p_{n,\lambda}^2 + \omega_n^2 q_{n,\lambda}^2 \Big), \label{eq:HF-sho}
% \end{equation}
% where $q_{n,\lambda}(t)$ is the canonical coordinate and
% \(
% p_{n,\lambda}(t) = \dot q_{n,\lambda}(t)
% \)
% is its conjugate momentum. Thus every cavity mode $(n,\lambda)$ is exactly a classical harmonic oscillator of unit mass. Quantization of the EM field follows by promoting the $(q,p)$ variables to operators that follow commutation relations, $[\hat q_{n,\lambda},\hat p_{m,\lambda'}]=i\hbar\,\delta_{nm}\delta_{\lambda\lambda'}$. 

% The next step is to perform a series of unitary operations onto the Hamiltonian. 
% The goal is to move the coupling from the matter degrees of freedom to the light degrees of freedom.  
% Not only this will bring us closer to Eq.~REF, but it also allows us to interpret the coupling as the additional polarization of the EM field or, within the SHO picture, the finite displacement of oscilattor positions due to finite polarizations. 
% To this end, one important simplification we can make is that we are chiefly interested in the liquid confined at the center of the cavity. 
% Due to the separation of lengthscales, we can only write the vector potential evaluated at the center. Since the mode functions are equal to $\sqrt{2/V}$, we obtain
% \begin{equation}
% \bm A(\bm r) \approx \bm A(0) = \sum_{n,\lambda} \sqrt{\frac{2}{V}} q_{n,\lambda} \hat{\bm e}_\alpha \,,
% \end{equation}
% In general cavities, the factor $\sqrt{2/V}$ is often replaced with a generic constant $\lambda$, which can vary further depending on the cavity geonetry and boundary conditions. From hereon, we replace $\sqrt{2/V}$ with $\lambda$. Keep in mind the scaling $\lambda \sim 1/\sqrt{V}$, since molecules exist at constant density $\rho$, implies that $\lambda \sim 1/\sqrt{N}$ where $N$ is number of particles. This is a classic result that implies that the coupling effect of cavity modes vanish in the thermodynamic limit. However, we shall see that such argument does not discount what occurs when the systme is at non-equilibrium. 

% Next, being a linear combination of all canonical positions, we can construct the following unitary operation that achieves a momentum boost that eliminates the interaction term, but moves them over to the momentum of the oscillators. In particular, noting that the interaction term boils down to $Z_j \bm A(0)$, we dot product this with matter position and thus defining the following unitary
% \begin{equation}
% \hat{U}_L = \exp \left[-\frac{\mathrm{i}}{\hbar}\bm \mu \cdot \bm A(0) \right] = \exp \left[-\frac{\mathrm{i}}{\hbar} \sum_{n,\lambda} \lambda (\bm \mu \cdot \bm{ \hat{e}}_{\lambda}) q_{n,\lambda} \right] 
% \end{equation}
% To see, how we can apply this to transform the system we note its action on the relevant operators is (Baker–Campbell–Hausdorff; using $[\hat q,\hat p]=i\hbar$ and $[\hat{\boldsymbol\mu},\hat q]=0$)
% \begin{align}
% \hat U_{\!L}^{\dagger}\, \hat p_{n,\lambda}\, \hat U_{\!L} &= \hat p_{n,\lambda} + \lambda \bm{ \hat{e}}_{\lambda}\!\cdot\!\hat{\boldsymbol\mu}, \label{eq:p-shift}\\
% \hat U_{\!L}^{\dagger}\, \hat q_{n,\lambda}\, \hat U_{\!L} &= \hat q_{n,\lambda}, \\
% \hat U_{\!L}^{\dagger}\, \hat{\mathbf p}_j\, \hat U_{\!L} &= \hat{\mathbf p}_j + Z_j\, \mathbf{A}(0), \\
% \hat U_{\!L}^{\dagger}\, \mathbf{A}(0)\, \hat U_{\!L} &= \mathbf{A}(0),
% \end{align}
% so that \emph{the minimal-coupling momentum} transforms as
% \begin{equation}
% \hat U_{\!L}^{\dagger}\, \big(\hat{\mathbf p}_j - Z_j \mathbf{A}(0)\big)\, \hat U_{\!L} = \hat{\mathbf p}_j .
% \end{equation}
% Thus, applying $\hat U_{\!L}$ to $\hat H$ thus eliminates the vector potential from the matter kinetic energy and shifts the \emph{light} momenta in $\hat H_\mathrm{EM}$:
% \begin{align}
% \hat H' \equiv \hat U_{\!L}^{\dagger} \hat H \hat U_{\!L}
% &= \underbrace{\left[\sum_i \frac{\hat{\mathbf P}_i^{\,2}}{2M_i} + \sum_i \frac{\hat{\mathbf p}_i^{\,2}}{2m_i} + \hat V_{nn}+\hat V_{ee}+\hat V_{ne}\right]}_{\displaystyle \hat H_\mathrm{M}\ \text{(no vector potential, } \bm A\text{)}} \nonumber\\
% &\quad + \frac12 \sum_{n,\lambda}\bigg[\Big(\hat p_{n,\lambda} + \lambda \bm{\hat{e}}_{\lambda}\!\cdot\!\hat{\boldsymbol\mu}\Big)^{\!2} + \omega_n^2 \hat q_{n,\lambda}^2\bigg].
% \label{eq:HL-dE}
% \end{align}
% The above Hamiltonian is also known as the \emph{length-gauge} ($d\!\cdot\!E$) Hamiltonian. Next, we can introduce the $90^\circ$ light phase rotation which transforms the light harmonic oscillator degrees of freedom into:
% \begin{gather}
% \hat U_{\pi/2} \equiv \exp\!\left[i\frac{\pi}{2}\sum_{n,\lambda}\hat a_{n,\lambda}^\dagger \hat a_{n,\lambda}\right] \,,
% \\
% \hat U_{\pi/2}^{\dagger}\,\hat q_{n,\lambda}\,\hat U_{\pi/2} = -\,\frac{\hat p_{n,\lambda}}{\omega_n}  \,,
% \\
% \hat U_{\pi/2}^{\dagger}\,\hat p_{n,\lambda}\,\hat U_{\pi/2} = \omega_n \hat q_{n,\lambda} \,.
% \end{gather}
% Acting with $\hat U_{\pi/2}$ on \eqref{eq:HL-dE} and renaming the rotated variables back to $(\hat q,\hat p)$ for notational simplicity leads to
% \begin{align}
% \hat H_{\mathrm{PF}}
% &= \hat U_{\pi/2}^{\dagger}\, \hat H_{d\cdot E}\, \hat U_{\pi/2} \nonumber\\
% &= \hat H_\mathrm{M}
% + \frac12 \sum_{n,\lambda}\left[
% \hat p_{n,\lambda}^2 + \omega_n^2\!\left(\hat q_{n,\lambda} - \frac{\lambda \bm{\hat{e}}_{\lambda}\!\cdot\!\hat{\boldsymbol\mu}}{\omega_n}\right)^{\!2}
% \right] .
% \label{eq:PF-final}
% \end{align}
% Equation~\eqref{eq:PF-final} is the multi-mode Pauli–Fierz Hamiltonian in “coordinate-displacement” form. 
% It is exactly of the Eq.~(56) type: the \emph{light coordinate} is displaced by a \emph{dipolar amount}. 
% In light of how we write everything, we can  re-write the interaction term as
% \begin{gather}
% \hat{H} = \hat{H}_\mathrm{M}+\hat{H}_\mathrm{F} \label{eq:H}
% \\
% \hat{H}_\mathrm{M} = \sum_i \frac{\hat{\mathbf P}_i^{\,2}}{2M_i} + \sum_i \frac{\hat{\mathbf p}_i^{\,2}}{2m_i} + \hat V_{nn}+\hat V_{ee}+\hat V_{ne} \label{eq:HM-final}
% \\
% \hat{H}_\mathrm{F} = \frac12 \sum_{n,\lambda}\left[
% \hat p_{n,\lambda}^2 + \omega_n^2\!\left(\hat q_{n,\lambda} - \frac{\lambda \hat{\mu}_\alpha }{\omega_n}\right)^{\!2}
% \right] \label{eq:HF-final}
% \end{gather}

% At this stage, we note several things: nuclei move more slowly than the the electrons but the EM resonance frequency is set to nuclei vibrational motion, hence operating at the same timesacles. Thus, we may employ the BO +adiabatic approximation in which we treat the electronic states separately from the nuclei+light DOFs. The electronic structure problem is 
% \begin{gather}
% \hat{H}_\mathrm{e} | \psi_g \rangle = V_\mathrm{BO}(\{ \bm{R}_i,q_{n,\lambda} \}) | \psi_g \rangle
% \\
% \hat{H}_\mathrm{e} = \sum_i \frac{\hat{\mathbf p}_i^{\,2}}{2m_i}+\hat V_{nn}+\hat V_{ee}+\hat V_{ne}
% \end{gather}
% Since the light degrees of freedom are very off resonant from electronic motion, we can also assume that $V_\mathrm{BO}(\{ \bm{R}_i,q_{n,\lambda} \}) \approx V_\mathrm{BO}(\{ \bm{R}_i\})$. 
% Long story short, in the BO approximation, where we want to treat dynamics of the slower moving DOGfs classically, the potential is governed by a mean-field potential generated by the the electronic ground state. And so, the classical potential energy is
% \begin{gather}
% V(\{ \bm{R}_i,q_{n,\lambda} \}) = V_\mathrm{BO}(\{ \bm R_i \}) + V_\mathrm{F}(\{\bm R_i, q_{n,\lambda}\})
% \\
% V_\mathrm{F} = \sum_{n,\lambda} \frac{1}{2} \omega_n q_{n,\lambda}^2 -\lambda \langle \psi_g |\hat{\mu}_\alpha |\psi_g \rangle  q_{n,\lambda}+\frac{\lambda^2}{2 \omega_n} \langle \psi_g | \hat{\mu}_\alpha^2 |\psi_g \rangle
% \end{gather}
% The term $\langle \psi_g | \hat{\mu}_\alpha^2 |\psi_g \rangle$ refers to the electronic zero-point fluctuation of the dipole moment operator. This is a higher-order effect that we don't care about so we can employ a mean-field operation $\langle \psi_g | \hat{\mu}_\alpha^2 |\psi_g \rangle \approx \left(\langle \psi_g | \hat{\mu}_\alpha |\psi_g \rangle \right)^2$. Defining the classical dipole moment of the system as $\mu_\alpha$, we can then write the interactino back as a displaced lcassical harmonic oscillator, and thus:
% \begin{gather}
% H  = H_\mathrm{M}+H_\mathrm{EM} \label{eq:H}
% \\
% H_\mathrm{M} = \sum_i \frac{\mathbf P_i^{\,2}}{2M_i} + V_\mathrm{BO}(\{ \bm R_i \}) \label{eq:HM-final}
% \\
% H_\mathrm{EM} = \sum_{n,\lambda} \frac{p_{n,\lambda}^2}{2}+ \frac{1}{2} \omega_n \left(q_{n,\lambda} -\frac{\lambda \mu_\alpha(\{ \bm R_i \})}{\omega_n} \right)^2
% \end{gather}
% The Eq.~REFs provide a simple approach for investigating cavity effects classically in a molecular system. One simply adopts a model for the PES and dipole moment ($V_\mathrm{BO}(\{ \bm R_i\{), \mu_\alpha$). THis can be done through a ladder of increasingly higher accurate models: non-polarizable empirical force fields, polarizable many-body force fields, machine learning (ML) models (for both PES and dipole moment), and first principles calculations (either Hartree/DFT or higher level). In this work, we adopt the simplest  classical force field model that is non-polarizable---writing the partial charges of nuclei as $q_i$, we can write the dipole moment $\bar{\bm \mu} = \sum_{i=1}^{N_\mathrm{n}} q_i \bm R_i$. 

% At this stage, the further simplifications that we can make is the fact that we only retain the lowest mode of the EM field---an approximation well-justified by the separation of energy scales between the fastest most relevant molecular processes and the next highest mode of the cavity, this brings us to Eq.~REFs.


% % In this section, we provide a short derivation on obtaining the model Hamiltonian given. As literature may focus on electronic excitations and a quantum treatment of the problem, emphasis is given here on a fully classical treatment of the problem from the very beginning, which is more well suited from our problem. The presentation here dies straight into the heart of the problem. For more fundamental discourse on the topic, see REFs. 

% % We begin with matter as a colelction of negatively charged electrons and positively charged nuclei. Within the Born-Oppenheimer approximation, the nuclei moves along the ground-state potential energy surface (PES) generated by the ground-state electronic configuration. This potential energy can be denoted as $V(\{ \bm R_i\})$. The nuclei, treated as classically, can be computed as follows:
% % \begin{equation}
% % H_\mathrm{M} = \sum_i \frac{| \bm P_i|^2}{2 M_j} + V(\{ \bm R_i \})
% % \end{equation}

% % Next, we model the Hamiltonian of an electromagnetic field in a Fabry-Perot cavity. The simplest approach is to note that the solutions of Maxwell's equations in such a cavity is that of a purely transverse electromagnetic field, whose wave vector $\bm k$ is perpendicular to its mirrors. Denoting the polarization directions as $\{ \bm e_\alpha \}$, we can write the electric and magnetic vector fields over space $\bm r$ separately as $ \bm E_{\perp}(\bm r)$ and $\bm B_{\perp}(\bm r)$ and the Hamiltonian is 
% % \begin{equation}
% % H_\mathrm{C} = \frac{\epsilon_0}{2} \int d^3 \bm r \left[ | \bm E_{\perp}(\bm r)|^2 +c^2 |\bm B (\bm r)|^2 \right] \,.
% % \end{equation}
% % In the most general treatment of EM theory, both fields are typically modeled by a set of scalar and vector potentials $(\phi, \bm A)$. However, the purely transverse nature of the fields and coupled with the use of Coulomb gauge implies that $\phi = 0$ everywhere, and that $\bm A$ is also purely transverse. From here, we can work in terms of the normal modes of the vector potential, satisfying the appropriate boundary conditions of the cavity as follows:
% % \begin{equation}
% % \bm A(\bm r) = \sum_{n \lambda} \frac{\hat{\bm e}_{\lambda}}{\omega_{n}} \sqrt{\frac{\hbar \omega_{n}}{2 \epsilon_0 V}} \left(a_{n  \lambda} e^{i n \hat{\bm k} \cdot \bm r}+a_{n \lambda }^\dagger e^{-i n \hat{\bm k} \cdot \bm r} \right)
% % \end{equation}
% % And the electric ang magnetic fields are  


% % \begin{enumerate}
% %     \item Note that we should begin with saying "Oh here is the Pauli-Fierz Hamiltonian" and we have such and such potential energy and interactions and that we are coupled to a classical model of atoms with fixed charges, but do mention that such model is meant to approximate the real ground-state Born-Oppenheimer PES. 
% %     \item Begin applying the single-mode approximation and then the point-dipole approximation. At some point, we need to do gauge transformation and a canonical transformation so that we bring it to a proper harmonic oscillator form. 
% % \end{enumerate}

% \section{Modeling the Baths}

% From Eq.~REFs, classical equations of motion can be generated via Hamilton's EOMs. But this is only part of the story; for a system at thermal equilibrium, which is then pushed away to non-equilibirium states (as we shall do in this work). We need to model the light-matter system as an open system, exchanging energy with the environment at rates realistic to the given study. In our setup, there are two main types of baths:
% \begin{enumerate}
%     \item \textit{Cavity bath:} Even outside of 

%     \begin{gather}
%     d q_{\lambda} = p_\alpha dt
%     \\
%     d p_\alpha = -\omega_0 (q_\alpha -\lambda \mu_\alpha/\omega_0) dt -\gamma_\mathrm{C} dt +\sqrt{2 k_\mathrm{B} T }dWx3(t)
%     \end{gather}

%     \item The
% \end{enumerate}
% The modeling of the molecular baths present the biggest challenge. Experimentally, we know that baths for such a macroscopic sample plays little role in structural evolution of the system but it does need to be efficient enough to, the latter implies fairly strong coupling to dissipate excess energy at a molecular timescale. Unfortunately, a corresponding Langevin bath to the nuclei DOFs can cause severe distortion to the original dynamics REFs. THis problem can be solved by adopting a bath model that corrects the microscopic dynamics minimally while still exchanging energy at the same timescale as a Langevin bath would. SUch a bath is known as the Bussi bath. 

% To derive this, we write the nuclei equations of motion with an additional correction for ce $\bm G_i$ that is left undetermined:
% \begin{gather}
% \dot{\bm R}_i = \frac{\bm{P_i}}{M_i}
% \\
% \dot{\bm P_i} = -\nabla_{\bm R_i} V_\mathrm{BO}+\bm G_i
% \end{gather}
% We note that the original dynamics is bath-free, which implies that $\dot{H}_\mathrm{M} = 0$. For a fixed magnitude of the correction force $G = | \bm G_i|$, the form of the correction force that minimizes the disturbance to the original dynamics one that is directed along the momentum direction. In other words,  
% \begin{equation}
% \bm G_i  = g \tilde{\bm P_i}
% \end{equation}
% where $\tilde{\bm P_i} = \bm P_i/|\bm P_i|$. To determine $g$, we simply note that $g$ needs to equate the energy exchange to the Langevin one. Next, we note down that the total energy, i.e., Hamiltonian, of the system exchanges energy with a Langevin bath as follows:
% \begin{equation}
% d H_\mathrm{M} = - \frac{K(t)-\bar{K}}{\tau} dt+2\sqrt{\frac{\bar{K} K(t)}{3 N_\mathrm{n} \tau}} d W(t)
% \end{equation}
% where $\eta(t)$ is white noise and $3 N_\mathrm{n}$ refers the total number of positional degrees of freedom we have. Since time evolution of total energy only depends on the instantaneous total kinetic energy $K(t)$, the most general form of the correction force must take the following form:
% \begin{equation}
% \bm G_i dt = \bm P_i \left[ -\gamma (K) dt +\sqrt{D(K)} d W(t)\right]
% \end{equation}
% And by Ito's lemma, the change in total energy using the EOMs in REF is:
% \begin{equation}
% \dot{H}_\mathrm{M} = \left(-2\gamma +D \right) K(t)+ 2\sqrt{D}K(t) dW(t)
% \end{equation}
% This implies that the effective friction and diffusion constant is
% \begin{gather}
% \gamma(K) = \frac{1}{2 \tau} \left(1-\left(1-\frac{1}{3 N_\mathrm{n} }\right) \frac{\bar{K}}{K} \right)
% \\
% D(K) = \frac{\bar{K}}{3 N_f \tau K}
% \end{gather}

\section{Theory}\label{app:theory}


\subsection{Deriving the Light-Matter Hamiltonian}\label{app:hamiltonian}

In this section, we present a self-contained derivation of the light-matter Hamiltonian employed in this study, with particular emphasis on obtaining a classical treatment suitable for molecular dynamics (MD) simulations.
We begin with the minimal coupling Hamiltonian $\hat{H}$ describing the interaction between matter---composed of negatively charged electrons and positively charged nuclei---and an electromagnetic (EM) field:
\begin{gather}
\hat{H} = \hat{H}_\mathrm{M}+\hat{H}_\mathrm{EM} \label{eq:H}
\\
\hat{H}_\mathrm{M} = \sum_{i=1}^{N_\mathrm{n}} \frac{| \hat{\bm P}_i -Z_i e \hat{\bm A}(\bm R_i)|^2}{2 M_i} + \sum_{j=1}^{N_e} \frac{| \hat{\bm p}_j +e \hat{\bm A}(\bm r_j)|^2}{2 m_j} \nonumber
\\
+ \hat{V}_\mathrm{nn}+ \hat{V}_\mathrm{ee}+ \hat{V}_\mathrm{ne} \label{eq:HM}
\\
\hat{H}_\mathrm{EM} = \frac{\epsilon_0}{2}\int_\Omega d^3 \bm r  \ \big( | \hat{\bm E}(\bm r)|^2+c^2 | \hat{\bm  B}(\bm r)|^2 \big) \label{eq:HF}
\end{gather}
Here, $\hat{H}_\mathrm{F}$ and $\hat{H}_\mathrm{M}$ represent the electromagnetic field and matter components of the Hamiltonian, respectively.
The first two terms in $\hat{H}_\mathrm{M}$ describe the kinetic energy of $N_\mathrm{n}$ nuclei (charge $+Z_i e$, mass $M_i$) and $N_e$ electrons (charge $-e$, mass $m_j$) coupled to the vector potential $\hat{\bm A}$ of the EM field.
The remaining terms $\hat{V}_\mathrm{nn}$, $\hat{V}_\mathrm{ee}$, and $\hat{V}_\mathrm{ne}$ represent nucleus-nucleus, electron-electron, and nucleus-electron Coulombic interactions, respectively.
In Eqs.~\eqref{eq:H}--\eqref{eq:HF}, we adopt the Coulomb gauge where $\nabla \cdot \hat{\bm A} = 0$ and incorporate the scalar potential $\phi$ into the static Coulombic interactions between charged particles.
Consequently, the EM fields are purely transverse.

Next, we transform the field Hamiltonian into a harmonic oscillator form. 
To this end, we write the EM energy density in terms of the vector potential, using the relations $\hat{\bm E} = -\partial_t \hat{\bm A}$ and $\hat{\bm B} = \nabla \times \hat{\bm A}$ to obtain:
\begin{equation}
\hat{H}_\mathrm{F}  = \frac{\epsilon_0}{2} \int_\Omega d^3 \bm r \ \left(\left|\partial_t \hat{\bm A} \right|^2
+ c^2 |\nabla \times \hat{\bm A}|^2 \right) \,.
\end{equation}
Next, we consider a one-dimensional Fabry-Pérot cavity with mirrors at $z=\pm L/2$ and expand the transverse vector potential in terms of normal mode functions:
\begin{equation}
\hat{\bm A}(\bm r,t) = \sum_{n,\alpha} \hat{q}_{n,\alpha}(t)\, \bm u_{n,\alpha}(\bm r),
\quad
\bm u_{n,\alpha}(\bm r)= u_n(z)\,\hat{\bm e}_\alpha,
\end{equation}
where $\hat{\bm e}_\alpha \perp \hat{\bm z}$ are fixed polarization vectors and
\begin{equation}
u_n(z) = \sqrt{\frac{2}{V}}\,\sin \left(\frac{n\pi}{L}\left(z+\frac{L}{2}\right)\right),
\quad \omega_n=\frac{n\pi c}{L}
\end{equation}
with cavity volume $V=A L$ and cross-sectional area $A$.
The mode functions satisfy orthogonality:
\begin{equation}
\int_V d^3\bm r\, \bm u_{n,\alpha} \cdot \bm u_{m,\alpha'} = \delta_{nm}\delta_{\alpha\alpha'} \,.
\end{equation}
Substituting this expansion into the field Hamiltonian and utilizing orthogonality causes each mode to decouple, yielding:
\begin{equation}
\hat{H}_\mathrm{F} = \frac{1}{2} \sum_{n,\alpha} \Big( \hat{p}_{n,\alpha}^2 + \omega_n^2 \hat{q}_{n,\alpha}^2 \Big) \label{eq:HF-sho}
\end{equation}
where $\hat{q}_{n,\alpha}(t)$ is the canonical coordinate and $\hat{p}_{n,\alpha}(t)$ is its conjugate momentum.
Each cavity mode $(n,\alpha)$ thus behaves as a quantum harmonic oscillator with unit mass.

A crucial simplification arises when molecules are confined at the cavity center.
Due to the separation between molecular and cavity length scales, we can approximate the vector potential as spatially uniform.
Since the mode functions equal $\sqrt{2/V}$ at the cavity center, we obtain:
\begin{equation}
\hat{\bm A}(\bm r) \approx \hat{\bm A}(0) = \sum_{n,\alpha} \sqrt{\frac{2}{V}} \hat{q}_{n,\alpha} \hat{\bm e}_\alpha \,.
\end{equation}
In general cavity geometries, the factor $\sqrt{2/V}$ is often replaced with a generic coupling constant $\lambda_\alpha$, which varies depending on geometry and boundary conditions.
Henceforth, we replace $\sqrt{2/V}$ with $\lambda_\alpha$, noting the scaling $\alpha \sim 1/\sqrt{V}$.
Since molecules exist at constant density $\rho$, this implies $\lambda\sim 1/\sqrt{N}$ where $N$ is the number of particles.
Thus, in the thermodynamic limit, a polariton formation arises even when the single-molecule coupling constant approaches zero. 
%This result indicates that coupling effects of individual cavity modes vanish in the thermodynamic limit, though this does not preclude significant effects when the system is driven to an initial non-equilibrium state.

The next step in our derivation involves unitary transformations that transfer the coupling from matter to EM field degrees of freedom.
This approach not only simplifies the Hamiltonian but also allows interpretation of the coupling as additional polarization of the EM field or, within the harmonic oscillator picture, as a finite displacement of equilibrium oscillator positions. 

The first transformation eliminates the interaction term from the kinetic energy of matter by transferring it to the light momenta.
Since the interaction term involves the total molecular dipole moment $\hat{\bm \mu} = \sum_i Z_i e \hat{\bm R}_i - \sum_j e \hat{\bm r}_j$, we define:
\begin{equation}
\hat{U}_L = \exp \left[-\frac{i}{\hbar}\hat{\bm \mu} \cdot \hat{\bm A}(0) \right] = \exp \left[-\frac{i}{\hbar} \sum_{n,\alpha} \alpha \mu_\alpha \hat{q}_{n,\alpha} \right] 
\end{equation}
where $\mu_\alpha = \bm\mu \cdot \hat{\bm e}_\alpha$. The action of this transformation on the relevant operators follows from the Baker-Campbell-Hausdorff formula.
Using the commutation relations $[\hat{q}_{n,\alpha},\hat{p}_{n^\prime,\alpha^\prime}]=i\hbar \delta_{nn^\prime} \delta_{\alpha\alpha^\prime}$ and $[\hat{\bm\mu}_\alpha,\hat{q}_{n,\alpha^\prime}]=0$, we obtain:
\begin{align}
\hat{U}_{L}^{\dagger}\, \hat{p}_{n,\alpha}\, \hat{U}_{L} &= \hat{p}_{n,\alpha} + \alpha \hat{\bm e}_{\alpha} \cdot \hat{\bm\mu} \label{eq:p-shift}\\
\hat{U}_{L}^{\dagger}\, \hat{q}_{n,\alpha}\, \hat{U}_{L} &= \hat{q}_{n,\alpha} \\
\hat{U}_{L}^{\dagger}\, \hat{\bm P}_i\, \hat{U}_{L} &= \hat{\bm P}_i - Z_i e\, \hat{\bm A}(0) \\
\hat{U}_{L}^{\dagger}\, \hat{\bm p}_j\, \hat{U}_{L} &= \hat{\bm p}_j + e\, \hat{\bm A}(0) \\
\hat{U}_{L}^{\dagger}\, \hat{\bm A}(0)\, \hat{U}_{L} &= \hat{\bm A}(0) \,.
\end{align}
Consequently, the minimal-coupling momenta transform as:
\begin{align}
\hat{U}_{L}^{\dagger}\, \big(\hat{\bm P}_i + Z_i e \hat{\bm A}(0)\big)\, \hat{U}_{L} &= \hat{\bm P}_i \\
\hat{U}_{L}^{\dagger}\, \big(\hat{\bm p}_j - e \hat{\bm A}(0)\big)\, \hat{U}_{L} &= \hat{\bm p}_j \,.
\end{align}
Applying $\hat{U}_{L}$ to the full Hamiltonian eliminates the vector potential from the matter kinetic energy and shifts the light momenta in $\hat{H}_\mathrm{F}$:
\begin{align}
\hat{H}' \equiv \hat{U}_{L}^{\dagger} \hat{H} \hat{U}_{L}
&= \underbrace{\left[\sum_i \frac{\hat{\bm P}_i^{\,2}}{2M_i} + \sum_j \frac{\hat{\bm p}_j^{\,2}}{2m_j} + \hat{V}_{nn}+\hat{V}_{ee}+\hat{V}_{ne}\right]}_{\displaystyle \hat{H}_\mathrm{M}\ \text{(no vector potential)}} \nonumber\\
&\quad + \frac{1}{2} \sum_{n,\alpha}\bigg[\Big(\hat{p}_{n,\alpha} + \alpha \hat{\bm e}_{\alpha} \cdot \hat{\bm\mu}\Big)^{\!2} + \omega_n^2 \hat{q}_{n,\alpha}^2\bigg] \,,
\label{eq:HL-dE}
\end{align}
This is the length-gauge ($d \cdot E$) Hamiltonian.

We now apply a $90^\circ$ phase rotation that transforms the light harmonic oscillator degrees of freedom:
\begin{gather}
\hat{U}_{\pi/2} \equiv \exp\!\left[i\frac{\pi}{2}\sum_{n,\alpha}\hat{a}_{n,\alpha}^\dagger \hat{a}_{n,\alpha}\right] \\
\hat{U}_{\pi/2}^{\dagger}\,\hat{q}_{n,\alpha}\,\hat{U}_{\pi/2} = -\,\frac{\hat{p}_{n,\alpha}}{\omega_n} \\
\hat{U}_{\pi/2}^{\dagger}\,\hat{p}_{n,\alpha}\,\hat{U}_{\pi/2} = \omega_n \hat{q}_{n,\alpha} \,.
\end{gather}
Acting with $\hat{U}_{\pi/2}$ on Eq.~\eqref{eq:HL-dE} and relabeling the rotated variables as $(\hat{q},\hat{p})$ for notational simplicity yields:
\begin{align}
\hat{H}_{\mathrm{PF}}
&= \hat{U}_{\pi/2}^{\dagger}\, \hat{H}'\, \hat{U}_{\pi/2} \nonumber\\
&= \hat{H}_\mathrm{M}
+ \frac{1}{2} \sum_{n,\alpha}\left[
\hat{p}_{n,\alpha}^2 + \omega_n^2\!\left(\hat{q}_{n,\alpha} + \frac{\lambda_{\alpha} \hat{\mu}_\alpha}{\omega_n}\right)^{\!2}
\right]
\label{eq:PF-final}
\end{align}
Equation~\eqref{eq:PF-final} is the multi-mode Pauli-Fierz Hamiltonian in coordinate-displacement form, where the light coordinate is displaced by an amount proportional to the molecular dipole moment.
We rewrite this compactly as:
\begin{gather}
\hat{H}_\mathrm{PF} = \hat{H}_\mathrm{M}+\hat{H}_\mathrm{F} \label{eq:H-final}
\\
\hat{H}_\mathrm{M} = \sum_i \frac{\hat{\bm P}_i^{\,2}}{2M_i} + \sum_j \frac{\hat{\bm p}_j^{\,2}}{2m_j} + \hat{V}_{nn}+\hat{V}_{ee}+\hat{V}_{ne} \label{eq:HM-final}
\\
\hat{H}_\mathrm{F} = \frac{1}{2} \sum_{n,\alpha}\left[
\hat{p}_{n,\alpha}^2 + \omega_n^2\!\left(\hat{q}_{n,\alpha} + \frac{\lambda_{\alpha} \hat{\mu}_\alpha }{\omega_n}\right)^{\!2}
\right] \label{eq:HF-final}
\end{gather}

At this stage, we exploit timescale separations in the system.
Nuclei move more slowly than electrons, but the EM resonance frequency typically matches nuclear vibrational motion, placing both on similar timescales.
We therefore employ the Born-Oppenheimer (BO) approximation combined with an adiabatic treatment, solving the electronic structure problem separately from nuclear and electromagnetic degrees of freedom:
\begin{gather}
\hat{H}_\mathrm{e} | \psi_g \rangle = V_\mathrm{BO}(\{ \bm{R}_i,q_{n,\alpha} \}) | \psi_g \rangle \\
\hat{H}_\mathrm{e} = \sum_j \frac{\hat{\bm p}_j^{\,2}}{2m_j}+\hat{V}_{nn}+\hat{V}_{ee}+\hat{V}_{ne}
\end{gather}

Since electromagnetic modes are typically far off-resonant from electronic transitions, we approximate $V_\mathrm{BO}(\{ \bm{R}_i,q_{n,\alpha} \}) \approx V_\mathrm{BO}(\{ \bm{R}_i\})$.
Within the BO approximation, treating the slower degrees of freedom classically, the potential energy becomes:
\begin{gather}
V(\{ \bm{R}_i,q_{n,\alpha} \}) = V_\mathrm{BO}(\{ \bm{R}_i \}) + V_\mathrm{F}(\{\bm{R}_i, q_{n,\alpha}\}) \\
V_\mathrm{F} = \sum_{n,\alpha} \left[ \frac{1}{2} \omega_n^2 q_{n,\alpha}^2 +\lambda_\alpha \langle \psi_g |\hat{\mu}_\alpha |\psi_g \rangle q_{n,\alpha}+\frac{\lambda_\alpha^2}{2 \omega_n} \langle \psi_g | \hat{\mu}_\alpha^2 |\psi_g \rangle \right] \,.
\end{gather}
The term $\langle \psi_g | \hat{\mu}_\alpha^2 |\psi_g \rangle$ represents quantum fluctuations of the dipole moment operator=.
For the classical treatment, we employ a mean-field approximation: $\langle \psi_g | \hat{\mu}_\alpha^2 |\psi_g \rangle \approx \left(\langle \psi_g | \hat{\mu}_\alpha |\psi_g \rangle \right)^2$.
Defining the classical dipole moment as $\mu_\alpha(\{\bm{R}_i\}) = \langle \psi_g | \hat{\mu}_\alpha |\psi_g \rangle$, we obtain the classical light-matter Hamiltonian:
\begin{gather}
H  = H_\mathrm{M}+H_\mathrm{EM} \label{eq:H-classical}
\\
H_\mathrm{M} = \sum_i \frac{\bm{P}_i^{\,2}}{2M_i} + V_\mathrm{BO}(\{ \bm{R}_i \}) \label{eq:HM-classical}
\\
H_\mathrm{EM} = \sum_{n,\alpha} \left[ \frac{p_{n,\alpha}^2}{2}+ \frac{1}{2} \omega_n^2 \left(q_{n,\alpha} +\frac{\alpha \mu_\alpha(\{ \bm{R}_i \})}{\omega_n} \right)^2 \right] \label{eq:HF-classical}
\end{gather}

Equations~\eqref{eq:H-classical}--\eqref{eq:HF-classical} provide a practical framework for investigating cavity effects through classical molecular dynamics.
One can adopt models for the potential energy surface $V_\mathrm{BO}(\{ \bm{R}_i\})$ and dipole moment $\mu_\alpha(\{\bm{R}_i\})$ of increasing sophistication: non-polarizable empirical force fields, polarizable many-body force fields, machine learning models, or first-principles calculations.
In this work, we employ a simple non-polarizable classical force field.
Writing the partial charges of the nuclei as $z_i$, the dipole moment becomes $\bar{\bm \mu} = \sum_{i=1}^{N_\mathrm{n}} z_ie \bm{R}_i$.

A final simplification retains only the lowest-frequency EM mode.
This is well-justified by the energy scale separation between relevant molecular vibrational processes and higher cavity modes, leading to the single-mode equations employed in this study.

\subsection{Modeling the Bath Environments}
\label{app:bath}

Classical equations of motion can be generated from Eqs.~\eqref{eq:H-classical}--\eqref{eq:HF-classical} via Hamilton's equations.
However, for a system initially at thermal equilibrium that is subsequently driven to non-equilibrium states, we must model the light-matter system as open, exchanging energy with environmental baths at rates realistic for the given physical system.
Our setup incorporates two distinct thermal baths with fundamentally different coupling mechanisms.

\paragraph{Cavity bath.}
Even in high-quality optical cavities, photon loss occurs through imperfect mirrors, absorption, and scattering.
This dissipation is modeled through coupling to a harmonic bath representing the electromagnetic environment outside the cavity.

To derive the cavity Langevin dynamics, we couple the cavity mode coordinate $q_\alpha$ (frequency $\omega_\mathrm{c}$, unit mass) linearly to a continuum of bath oscillators in the Caldeira-Leggett form:
\begin{equation}
H_\mathrm{b} = \sum_j \left[\frac{p_j^2}{2m_j} + \frac{1}{2}m_j \omega_j^2 \left(x_j - \frac{c_j}{m_j \omega_j^2} q_\alpha\right)^2 \right]
\end{equation}
where $(x_j, p_j)$ are bath oscillator coordinates with frequencies $\omega_j$ and coupling constants $c_j$. Integrating out the bath degrees of freedom yields the generalized Langevin equation
\begin{equation}
\ddot{q}_\alpha(t) + \omega_\mathrm{c}^2 q_\alpha(t) + \int_0^t \!\! dt' \, \Gamma(t-t') \, \dot{q}_\alpha(t') = \eta_\alpha(t)
\end{equation}
where the memory kernel and noise correlation are related through the spectral density
\begin{equation}
J(\omega) = \frac{\pi}{2} \sum_j \frac{c_j^2}{m_j \omega_j} \delta(\omega - \omega_j)
\end{equation}
with the equations,
\begin{gather}
\Gamma(t) = \frac{2}{\pi} \int_0^\infty \!\! d\omega \, \frac{J(\omega)}{\omega} \cos(\omega t) \,, 
\\
\quad
\langle \xi_\alpha(t) \xi_\alpha(t') \rangle = k_\mathrm{B} T \, \Gamma(t-t') \,.
\end{gather}
In the Ohmic, Markovian limit where $J(\omega) \approx \gamma_\mathrm{c} \omega$ with a high-frequency cutoff, the memory kernel becomes instantaneous, $\Gamma(t) \approx 2\gamma_\mathrm{c} \delta(t)$, yielding standard Langevin dynamics in Eq.~REF (the $\lambda=0$ case).
The cavity lifetime $\tau_\mathrm{c}$ relates to the damping rate through $\gamma_\mathrm{c} = 1/(2\tau_\mathrm{c})$.
When the cavity mode is displaced by the molecular dipole, $q_\alpha \to q_\alpha + \lambda \mu_\alpha / \omega_\mathrm{c}$, the damping and noise statistics remain unchanged while the deterministic force acquires the cavity-molecule coupling term.

\paragraph{Molecular bath.}
The molecular environment presents a greater modeling challenge.
Experimentally, such baths play a limited role in structural evolution but must be efficient enough to dissipate excess energy on molecular timescales.
A direct Langevin approach applied to nuclear degrees of freedom can cause severe distortion of the intrinsic dynamics---particularly problematic for glassy systems where microscopic dynamics encode essential structural information.

This problem is resolved by the Bussi thermostat, which maintains the same energy exchange rate as a Langevin bath while minimally perturbing the microscopic dynamics.
The key insight is to apply stochastic corrections along the momentum direction rather than as external random forces
\begin{gather}
\dot{\bm{R}}_i = \frac{\bm{P}_i}{M_i} \\
\dot{\bm{P}}_i = -\nabla_{\bm{R}_i} V + \bm G_i(t)
\end{gather}
For a system with instantaneous kinetic energy $K(t)$ and target kinetic energy $\bar{K} = \frac{3}{2}Nk_\mathrm{B} T$, the correction takes the form
\begin{equation}
\bm{G}_i = \bm{P}_i \left[ -\gamma(K) + \sqrt{2 D(K)} \eta_i(t) \right]
\end{equation}
where the generalized friction and diffusion coefficients are
\begin{equation}
\gamma(K) = \frac{1}{2 \tau_\mathrm{b}} \left(1-\left(1-\frac{1}{3N}\right) \frac{\bar{K}}{K} \right) \,, \quad
D(K) = \frac{\bar{K}}{3N \tau_\mathrm{b} K} \,.
\end{equation}
Here $\tau_\mathrm{b}$ is the molecular thermostat time constant.
These coefficients are chosen to reproduce the energy exchange statistics of a Langevin bath while preserving velocity autocorrelation functions at short times---critical for maintaining the correct microscopic dynamics in supercooled liquids.
The resulting thermostat samples the canonical ensemble exactly while introducing minimal artifacts in dynamical observables.



\section{Computational Details}

\subsection{Numerical Integration Schemes}

All simulations employ split-operator integration schemes that separate deterministic equations of motion from stochastic thermostatting.
For molecular degrees of freedom, we use a two-step velocity-Verlet propagator combined with the Bussi stochastic velocity-rescaling thermostat.
At each timestep $\Delta t$, the molecular coordinates $(\bm{R}_i, \bm{P}_i)$ for particle $i$ with mass $M_i$ evolve as follows.
The first half-step advances momenta by
\begin{equation}
\bm{P}_i\left(t + \frac{\Delta t}{2}\right) = \bm{P}_i(t) + \frac{1}{2} \bm{F}_i(t) \Delta t \,,
\end{equation}
where $\bm{F}_i(t) = -\partial V / \partial \bm{R}_i$ is the total force on particle $i$ arising from intermolecular, intramolecular, and cavity coupling interactions.
The momentum is then rescaled by a stochastic factor $\alpha$ computed from the current kinetic energy
\begin{equation}
K = \sum_{i=1}^{N} \frac{|\bm{P}_i(t + \Delta t/2)|^2}{2M_i} \,,
\end{equation}
the temperature $T$, the number of degrees of freedom $f$, and the thermostat time constant $\tau_\mathrm{b}$:
\begin{equation}
\alpha = \sqrt{\gamma + \frac{T}{2K}(1-\gamma)\left(R_\Gamma + R_\mathrm{N}^2\right) + 2R_\mathrm{N}\sqrt{\frac{T}{2K}(1-\gamma)\gamma}} \,,
\end{equation}
where $\gamma = \exp(-\Delta t/\tau_\mathrm{b})$, $R_\mathrm{N}$ is a standard normal random variate, and $R_\Gamma$ is a gamma-distributed random variate with shape parameter $(f-1)/2$.
After rescaling all momenta by $\bm{P}_i \to \alpha \bm{P}_i$, positions update according to
\begin{equation}
\bm{R}_i(t + \Delta t) = \bm{R}_i(t) + \frac{\bm{P}_i(t + \Delta t/2)}{M_i} \Delta t \,.
\end{equation}
Forces are then recalculated at the new positions $\bm{F}_i(t + \Delta t)$, after which the momentum undergoes a second rescaling by a freshly computed factor $\alpha'$ followed by the second half-step:
\begin{equation}
\bm{P}_i(t + \Delta t) = \alpha' \bm{P}_i\left(t + \frac{\Delta t}{2}\right) + \frac{1}{2} \bm{F}_i(t + \Delta t) \Delta t \,.
\end{equation}
This integration scheme exactly samples the canonical ensemble while preserving the velocity autocorrelation function at short times, making it ideal for glassy dynamics.
Translational and rotational degrees of freedom are rescaled independently using their respective kinetic energies.
Equipartition is verified throughout all simulations (\cref{fig:equipartition}), confirming correct thermalization.

The cavity mode follows Langevin dynamics integrated via a modified velocity-Verlet scheme.
For the cavity mode $\alpha$ with canonical coordinates $(q_\alpha, p_\alpha)$ and unit mass, the first half-step advances the momentum by
\begin{equation}
p_\alpha\left(t + \frac{\Delta t}{2}\right) = p_\alpha(t) + \frac{1}{2} f_\alpha(t) \Delta t \,,
\end{equation}
where $f_\alpha(t) = -\omega_\mathrm{c}^2(q_\alpha + \lambda \mu_\alpha / \omega_\mathrm{c})$ is the effective force combining the harmonic restoring potential and cavity-molecule coupling.
The position then updates according to
\begin{equation}
q_\alpha(t + \Delta t) = q_\alpha(t) + p_\alpha\left(t + \frac{\Delta t}{2}\right) \Delta t \,.
\end{equation}
In the second half-step, stochastic forces are added through a modified acceleration:
\begin{equation}
f_\alpha(t + \Delta t) \to f_\alpha(t + \Delta t) - \gamma_\mathrm{c} \, p_\alpha\left(t + \frac{\Delta t}{2}\right) + \xi_\alpha(t) \,,
\end{equation}
where $\gamma_\mathrm{c} = 1/(2\tau_\mathrm{c})$ is the friction coefficient and $\xi_\alpha(t)$ is Gaussian white noise with variance
\begin{equation}
\langle \xi_\alpha^2 \rangle = \frac{6 \gamma_\mathrm{c} k_\mathrm{B} T}{\Delta t}
\end{equation}
as required by the fluctuation-dissipation theorem.
The final momentum is then
\begin{equation}
p_\alpha(t + \Delta t) = p_\alpha\left(t + \frac{\Delta t}{2}\right) + \frac{1}{2} f_\alpha(t + \Delta t) \Delta t \,.
\end{equation}
This Langevin integrator maintains the cavity mode at temperature $T$ on timescale $\tau_\mathrm{c}$ while allowing efficient energy exchange with the electromagnetic environment.

\subsection{Adaptive Timestepping}

Simulating systems with disparate timescales---femtosecond Rabi oscillations, picosecond bath relaxation, and nanosecond structural dynamics---requires careful timestep control.
We implement adaptive timestepping where the integrator dynamically adjusts $\Delta t$ to maintain integration accuracy without sacrificing computational efficiency.
The algorithm estimates local truncation error by comparing the velocity-Verlet trajectory against a single explicit Euler step.
For a particle $i$ at position $\bm{R}_i(t)$ with momentum $\bm{P}_i(t)$ and mass $M_i$, define the velocity $\bm{v}_i(t) = \bm{P}_i(t)/M_i$ and acceleration $\bm{a}_i(t) = \bm{F}_i(t)/M_i$.
The velocity-Verlet predictor yields
\begin{equation}
\bm{R}_{i,\mathrm{VV}}(t + \Delta t) = \bm{R}_i(t) + \bm{v}_i(t) \Delta t + \frac{1}{2} \bm{a}_i(t) (\Delta t)^2 \,,
\end{equation}
while the explicit Euler method gives
\begin{equation}
\bm{R}_{i,\mathrm{E}}(t + \Delta t) = \bm{R}_i(t) + \bm{v}_i(t) \Delta t \,.
\end{equation}
The per-particle error estimate is
\begin{equation}
\delta_i = \left| \bm{R}_{i,\mathrm{VV}}(t + \Delta t) - \bm{R}_{i,\mathrm{E}}(t + \Delta t) \right| = \frac{1}{2} \left|\bm{a}_i(t)\right| (\Delta t)^2 \,,
\end{equation}
and the system-wide error is computed as the root-mean-square
\begin{equation}
\delta = \sqrt{\frac{1}{N} \sum_{i=1}^{N} \delta_i^2}
\end{equation}
over all $N$ particles.
The timestep is then adjusted according to
\begin{equation}
\Delta t' = \Delta t \sqrt{\frac{\varepsilon(t)}{\delta}} \,,
\end{equation}
where $\varepsilon(t)$ is the target error tolerance.
This approach automatically increases $\Delta t$ when forces are weak and decreases it during high-force events, maintaining accuracy while maximizing computational throughput.
Since this variable-timestep scheme is not symplectic, it trades long-term energy conservation for adaptability, making it well-suited for our constant-temperature simulations where thermostats regulate total energy.

To prevent numerical instabilities during protocol switches or when restarting from saved configurations, we employ error tolerance ramping.
The effective error tolerance evolves according to
\begin{equation}
\varepsilon(t) = \varepsilon_* - \left(\varepsilon_* - \varepsilon_0\right) e^{-t/\tau_*} \,,
\end{equation}
where
\begin{equation}
\varepsilon_0 = f_0 \varepsilon_*
\end{equation}
is the initial error tolerance set by the initial fraction $f_0$, and $\tau_*$ is the ramping time constant.
Starting with a conservative timestep at $t=0$, the integrator gradually increases $\Delta t$ as $\varepsilon(t)$ approaches $\varepsilon_*$.
In this work, we use
\begin{equation}
\varepsilon_* = 0.01 \,, \quad f_0 = 10^{-5} \,, \quad \tau_* = 50\text{ ps}
\end{equation}
for all production runs.
This choice balances numerical stability during the initial transient period against computational efficiency once the system equilibrates.
When $f_0 = 0$, ramping is disabled and the integrator uses fixed $\varepsilon_*$ throughout, which we employ only for well-equilibrated continuation runs.

\subsection{System Parameters and Force Field}

All simulations employ $N=250$ diatomic molecules ($500$ atoms total) in a cubic simulation box of size $L = 40$ with periodic boundary conditions, thus the system density is $\rho = 0.0078125$.
The molecular composition follows a $4:1$ mixture of $\mathrm{A}_2:\mathrm{B}_2$ molecules.
Each diatomic molecule carries a permanent dipole: one atom has charge $+0.3e$ and the other has charge $-0.3e$, with the charges oriented along the molecular bond axis.

The complete force field parameters are summarized in \cref{tab:forcefield}.
Harmonic bond parameters are chosen to yield vibrational frequencies of $\omega_{\mathrm{AA}} = 1560$ cm$^{-1}$ and $\omega_{\mathrm{BB}} = 2433$ cm$^{-1}$, matching O$_2$ and N$_2$ stretching modes, respectively.
Lennard-Jones parameters for $\mathrm{A}_2$ molecules are tuned to O$_2$, with remaining parameters scaled to preserve Kob-Andersen model characteristics.
The LJ potential is shifted to zero at the cutoff distance and truncated beyond $r_\mathrm{cut}$.
Coulombic interactions between partial charges are computed without truncation using the minimum image convention.
A cell-based neighbor list with buffer distance $\Delta r = 1.0$ Bohr is employed, with bonded pairs excluded from nonbonded interactions.

\begin{table}[t]
\caption{\label{tab:forcefield}Force field parameters for the molecular Kob-Andersen (mKA) model. Harmonic bond parameters: spring constant $k$ and equilibrium bond length $R^0$. Lennard-Jones parameters: well depth $\epsilon$ and size parameter $\sigma$. Cutoff radius $r_\mathrm{cut}$ applies to all LJ interactions. All quantities in atomic units (Hartree for energy, Bohr for length). Each diatomic molecule has opposing partial charges $\pm 0.3e$ on its two atoms, creating a permanent dipole.}
\begin{ruledtabular}
\begin{tabular}{lccc}
\multicolumn{4}{c}{\textbf{Intramolecular (Harmonic Bonds)}} \\
\hline
Bond type & $k$ (Hartree/Bohr$^2$) & $R^0$ (Bohr) & $\omega$ (cm$^{-1}$) \\
\hline
A--A & $0.73204$ & $2.2817$ & $1560$ \\
B--B & $1.4325$ & $2.0744$ & $2433$ \\
\hline
\multicolumn{4}{c}{} \\
\multicolumn{4}{c}{\textbf{Intermolecular (Lennard-Jones)}} \\
\hline
Pair type & $\epsilon$ (Hartree) & $\sigma$ (Bohr) & $r_\mathrm{cut}$ (Bohr) \\
\hline
A--A & $1.6685 \times 10^{-4}$ & $6.2304$ & $15.0$ \\
B--B & $8.3426 \times 10^{-5}$ & $5.4828$ & $15.0$ \\
A--B & $2.5028 \times 10^{-4}$ & $4.9832$ & $15.0$ \\
\hline
\multicolumn{4}{c}{} \\
\multicolumn{4}{c}{\textbf{Electrostatic (Coulombic)}} \\
\hline
\multicolumn{4}{l}{Partial charges: $\pm 0.3e$ per molecule} \\
\multicolumn{4}{l}{No truncation (minimum image convention)} \\
\end{tabular}
\end{ruledtabular}
\end{table}

\textbf{Thermostat parameters.}
Both molecular and cavity baths operate at temperature $T$ with time constants $\tau_\mathrm{b} = \tau_\mathrm{c} = 5.0$ ps for all production runs unless otherwise specified.
The Bussi thermostat rescales translational and rotational degrees of freedom independently to maintain proper equipartition.

\subsection{Simulation Protocol}

The nonthermal aging simulations reported in this work follow a three-stage protocol designed to generate statistically independent initial configurations and probe nonequilibrium dynamics following cavity activation.

\paragraph*{Initial configuration generation.}
We generate the initial molecular configuration by placing $250$ diatomic molecules on a simple cubic lattice at the target density.
Each dimer center is positioned at a lattice site, with the two atoms of each molecule separated by their equilibrium bond length $R^0$ along a randomly oriented direction.
Bond lengths are sampled from the equilibrium Boltzmann distribution for harmonic bonds at temperature $T$: $r \sim \mathcal{N}(R^0, \sqrt{k_\mathrm{B}T/k})$, where $k$ is the bond spring constant.
This initial configuration provides a reasonable starting point that avoids particle overlaps while respecting thermal fluctuations in intramolecular coordinates.

\paragraph*{Equilibration and sampling of independent configurations.}
Starting from the Stage 1 configuration, we perform a single long NVT equilibration run at $T = 100$ K without cavity coupling ($\lambda = 0$).
The molecular system is thermalized using the Bussi thermostat with a short time constant $\tau_\mathrm{b} = 1.0$ ps to accelerate equilibration.
From this trajectory, we extract $500$ configurations separated by the structural relaxation time $\tau_\mathrm{s} \approx 50$ ps at $T = 100$ K, ensuring statistical independence between samples.
These configurations serve as initial conditions for the $500$ independent production runs, enabling robust ensemble averaging of nonequilibrium observables.

\paragraph*{Production runs.}
Each of the $500$ configurations from Stage 2 serves as the initial condition for an independent production trajectory.
At $t = 0$, we introduce a cavity particle with position and momentum sampled from the thermal equilibrium distribution at $T = 100$ K, corresponding to a harmonic oscillator with frequency $\omega_\mathrm{c}$.
The cavity-molecule coupling is initially set to zero ($\lambda(t < t_0) = 0$), allowing the system to evolve under equilibrium conditions.
At $t_0 = 200$ ps, the coupling is instantaneously activated to the target value $\lambda(t \geq t_0) = \lambda_\mathrm{target}$ according to \cref{eq:varepsilon}, initiating the nonthermal aging process.
For $t > t_0$, adaptive timestepping is enabled with the parameters specified in the previous subsection, and both molecular ($\tau_\mathrm{b} = 5.0$ ps) and cavity ($\tau_\mathrm{c} = 5.0$ ps) thermostats maintain temperature $T$.
Each production run continues for a total simulation time sufficient to observe relaxation back toward equilibrium, typically $2$--$5$ ns depending on coupling strength.

\paragraph*{Observable sampling and output.}
Throughout the production runs, we monitor multiple observables at different frequencies.
Comprehensive energy decompositions (kinetic, potential, bond, Lennard-Jones, Coulombic, cavity) are recorded every $0.1$ ps to track energy flow between subsystems.
The intermediate scattering function $F_{\bm{k}}(t; t_\mathrm{w})$ is computed by averaging over $50$ wavevectors $\{\bm{k}_i\}$ uniformly distributed on a sphere of radius $k = 2\pi/L$, corresponding to the first shell in reciprocal space and approximately the first peak of the static structure factor $S_{\bm{k}}$.
Reference configurations for computing $F_{\bm{k}}(t; t_\mathrm{w})$ are saved every $200$ ps for waiting time analysis.
Full trajectory snapshots are written to disk every $50$ ps for post-processing and visualization.
All random number generators use replica-dependent seeds (seed = $10000 +$ replica) to ensure reproducibility while maintaining independence between trajectories.

\subsection{Implementation Details}

All simulations use the \texttt{cavHOOMD-blue} software package, our GPU-accelerated extension of the \texttt{HOOMD-blue} molecular dynamics framework.
The cavity-molecule coupling force is implemented in both CPU (C++) and GPU (CUDA) variants, with the GPU version utilizing kernel autotuning to optimize thread block configurations for the specific hardware.
Both implementations compute identical forces and maintain bit-level reproducibility when using the same random seed, differing only in performance characteristics.
Typical production runs on an NVIDIA A100 GPU achieve simulation rates of $10^6$ timesteps per hour for our $N=250$ molecular system.

\section{Constructing the Material Time}

Material-time translational invariance (MTTI), given by \cref{eq:MTTI}, allows us to reconstruct the material time $\xi(t)$ directly from ISF measurements at multiple waiting times. We invert \cref{eq:MTTI} to extract material time differences from measured correlations:
\begin{equation}
\xi(t+t_\mathrm{w})-\xi(t_\mathrm{w}) = \Phi^{-1}[\phi_{\bm{k}}(t; t_\mathrm{w})] \,.
\end{equation}
Since we define structural relaxation times through the criterion $\phi_{\bm{k}}(\tau_\mathrm{s}; t_\mathrm{w})=0.1$, we normalize $\Phi$ such that $\Phi^{-1}(0.1) = 1$. This normalization defines one unit of material time as the interval required for one structural relaxation event. For a sequence of waiting times $t_\mathrm{w}^{(i)} = i \Delta t_\mathrm{w}$ with $i=1,2,\ldots$, each measurement of $\tau_\mathrm{s}^{(i)}$ yields a constraint on the material time:
\begin{equation}
\xi\left(\tau_\mathrm{s}^{(i)}+t_\mathrm{w}^{(i)}\right)-\xi(t_\mathrm{w}^{(i)})=1 \,. \label{eq:xitime}
\end{equation}
\Cref{eq:xitime} provides a set of constraints that we solve recursively to reconstruct the full material time function $\xi(t)$.

To implement this procedure algorithmically, we reconstruct $\xi(t)$ from ISF measurements at $N_\mathrm{ref}$ independent references through the following steps:
\begin{enumerate}
\item For each reference $i$, measure the relaxation time $\tau_\mathrm{s}^{(i)}$ through linear interpolation between adjacent data points where $\phi_{\bm{k}}(t; t_\mathrm{w}^{(i)})$ crosses the criterion value. This defines the laboratory time pair $t_1^{(i)} = t_\mathrm{w}^{(i)} + \tau_\mathrm{s}^{(i)}$ and $t_2^{(i)} = t_\mathrm{w}^{(i)}$, which satisfy the constraint $\xi(t_1^{(i)}) - \xi(t_2^{(i)}) = 1$.
\item Starting from the initial condition $\xi(0)=0$, solve the constraint system sequentially. For each constraint pair $(t_2^{(i)}, t_1^{(i)})$, determine $\xi(t_2^{(i)})$ by interpolating from previously established material time values, then compute $\xi(t_1^{(i)}) = \xi(t_2^{(i)}) + 1$.
\item Construct a continuous material time function on a uniform grid through piecewise linear interpolation of all constraint points, extrapolating linearly beyond the last measured point using the slope determined from the final two points.
\item Enforce monotonicity of the ISF by requiring $\phi_{\bm{k}}(t_{j+1}; t_\mathrm{w}) \leq \phi_{\bm{k}}(t_j; t_\mathrm{w})$ for all adjacent time points through a forward-pass correction.
\end{enumerate}

To validate that the reconstructed material time correctly captures MTTI, we transform all ISF data using the common material time function and verify collapse onto a master curve. For each measurement $\phi_{\bm{k}}(t; t_\mathrm{w})$, we compute the corresponding material time difference $\Delta\xi = \xi(t+t_\mathrm{w}) - \xi(t_\mathrm{w})$ and plot $\phi_{\bm{k}}(t; t_\mathrm{w})$ versus $\Delta\xi$ for all waiting times. Successful collapse indicates that a single master function $\Phi_{\bm{k}}(\Delta\xi)$ describes the relaxation dynamics independent of waiting time and coupling strength (\cref{fig:fig4}b). 
We fit this master curve to the stretched exponential form given in the main text, which yields the universal stretching exponent $\beta = 0.55$ characterizing glassy relaxation in our system.



\end{document}
%
% ****** End of file apssamp.tex ******

